
\chapter{K3$\times$E geometry, Mukai dictionary, lattice polarization,
normal-ordered Gram cocycle, hybrid $\operatorname{Ran}^{\mathrm{hyb}}(E)$ carrier}
\label{ch:k3xe-setup}

The arithmetic source contains three named weight-five readings, and
their normalizations are different.  The automorphic chain produces the section
\[
\mathcal{D}_X = \Delta_5
\in M_5\!\bigl(\Sp_4(\Z),\,\nu_{\Delta_5}\bigr)
\qquad\text{(View \textcircled{1}; proved \cite{Borcherds95,GN}).}
\]
The Borcherds--Kac--Moody (BKM) chain produces the denominator
\[
\operatorname{den}(\mathfrak g_{\Delta_5}) = 64^{-1}\,\Delta_5(2Z)
\qquad\text{(View \textcircled{2}; proved \cite{GN}).}
\]
The compact Pfaffian chain produces, granted the retained D0-HN datum,
the retained quotient-orientation datum, the chosen Weyl lifts, and
the retained type-II wall-atlas datum of
Theorem~\ref{thm:G-four-obstruction-criterion} of
Appendix~\ref{app:obstruction-discharges}, the protected Pfaffian
\[
\operatorname{Pf}_{\mathrm{prot}}(\mathfrak D_X^{\mathrm{DI}}) = \mathcal D_X=\Delta_5
\qquad\text{(View \textcircled{3}; relative to the retained orientation data at the Pfaffian level).}
\]
The handle ``$\Delta_5$'' with no superscript is reserved for working
formulas; each occurrence names which of $\mathcal D_X$,
$\operatorname{den}(\mathfrak g_{\Delta_5})$, or
$\operatorname{Pf}_{\mathrm{prot}}(\mathfrak D_X^{\mathrm{DI}})$ is meant.
The cross-level identity \textcircled{2}$\leftrightarrow$\textcircled{3}
is a theorem at the Pfaffian and orientation level on \(K3\times E\),
granted \((\mathrm{D0})\), the retained quotient-orientation datum for
$(\mathrm{O1})$, the chosen Weyl lifts for $(\mathrm{O1}^+)$, and the
retained datum $(\mathrm{O2}_{\mathrm{atlas}})$
in Theorems~\ref{thm:G-D0}, \ref{thm:G-O1}, \ref{thm:G-O1-plus},
\ref{thm:G-O2-orbit-transport}.  The primitive Lie-superalgebra
recognition \(P_X^{\Pi,\mathrm{red}}\cong\mathfrak g_{\Delta_5}\)
requires the finite compact-source datum \((S1)\)--\((S10)\) of
Chapter~\ref{ch:pfaffian-dirac}.

\paragraph{The Mukai--Gram dictionary.}
The Mukai--Gram dictionary on \(K3\times E\) begins with the Mukai
lattice \(\widetilde H(S,\Z)\) of the K3 surface \(S\), of rank \(24\)
and signature \((4,20)\) \cite{Mukai1984}.  The even Kunneth sector
\[
\Gamma_X^{\mathrm{form}}
 =\widetilde H(S,\Z)\otimes H^{\mathrm{even}}(E,\Z)
 \cong \widetilde H(S,\Z)\oplus\widetilde H(S,\Z)
\]
is the formal \(K3\times E\) charge lattice used below; the two summands
are written \((Q,P)\).  For the Gram map this underlying Kunneth group
is read with the two componentwise K3 Mukai pairings; the tensor
Kunneth cross-pairing on \(H^{\mathrm{even}}(X,\Z)\) is a different
bilinear form.  The Gram-index lattice
\(\Gamma_{\mathrm{gram}}=\Z^3\) receives the quadratic shadow
\[
 \Pi_X(Q,P)=\bigl((Q,Q)/2,\,(Q,P),\,(P,P)/2\bigr).
\]
This map is not additive; its additivity defect is the symmetric
polarization \(B=-\delta\Pi_X\).  The additive lattice map is the
normal-ordered homomorphism
\[
 \overline\Pi_X:\widehat\Gamma_X\longrightarrow\Gamma_{\mathrm{gram}},
 \qquad
 \widehat\Gamma_X=\Gamma_X^{\mathrm{form}}\oplus_B\Gamma_{\mathrm{gram}}.
\]
After this normal ordering, \(\Gamma_{\mathrm{gram}}\) embeds into
\(\La^{2,1}_{II}\) and then into the type-IV ambient lattice
\(\La^{3,2}\) of Chapter~\ref{ch:singular-theta-lift}.  Detailed lattice
computations are in Appendix~\ref{appendix:mukai-gram-dictionary}; the
dictionary fixes the \(K3\times E\) charge convention.

The geometric and lattice data feed the three established
\(\Delta_5\) constructions.  The compact target is fixed at
$X = S \times E$ with $S$ a smooth complex
projective K3 surface and $E$ a smooth elliptic curve over $\C$; this is
a trivial-canonical threefold ($\dim_\C X = 3$, $K_X \cong \mathcal{O}_X$;
elementary from the product formula), so $n = \chi(\mathcal O_Y)$
appears directly.  The $(\kappa_{\mathrm{cat}},
\kappa_{\mathrm{ch}}^{\mathrm{Hodge}}, \kappa_{\mathrm{ch}}^{\mathrm{Heis}},
\kappa_{\mathrm{BKM}}, \kappa_{\mathrm{fiber}}) = (0, 0, 3, 5, 24)$ tuple
of \cite{CYQGVolIII} (the K3$\times$E row of the
Gritsenko--Cl\'ery catalogue) is the categorical-dimensional anchor;
every use of $\kappa$ carries one of these five subscripts.

The datum has a geometric part and an arithmetic part.  The geometric
part is the formal dyonic charge lattice
$\Gamma_X^{\mathrm{form}} = \Lambda_S \oplus \Lambda_S$ on the Mukai
sector, together with
$\Gamma_{X,\sigma}^{\mathrm{eff}} \subset \Gamma_X^{\mathrm{alg}}
\subset \Gamma_X^{\mathrm{form}}$.  The arithmetic part is the
quadratic Gram map
$\Pi_X : \Gamma_X^{\mathrm{form}} \to \Gamma_{\mathrm{gram}} = \Z^3$,
its symmetric bilinear polarization $B = -\delta\Pi_X$, the
normal-ordered extension
$\widehat\Gamma_X = \Gamma_X^{\mathrm{form}} \oplus_B
\Gamma_{\mathrm{gram}}$, and the additive map
$\overline\Pi_X(c, T) = \Pi_X(c) - T$.  The fixed-lift
raw-pushforward obstruction theorem forces the extension, and the
projection-locality obstruction forces the wrapped stratum in
$\operatorname{Ran}^{\mathrm{hyb}}(E) =
\operatorname{Ran}(E)_{b=0}\sqcup
\operatorname{Ran}^{\mathrm{wr}}(E)_{b>0}$.

\section{Smooth K3$\times$E threefold and Mukai lattice}
\label{sec:k3xe-mukai-lattice}

\subsection{The compact target}
\label{subsec:k3xe-compact-target}

\begin{definition}[K3$\times$E threefold; primitive geometric datum]
\label{def:k3xe-threefold}
A \emph{K3$\times$E threefold} is a triple $(X, S, E)$ in which $S$ is a
smooth complex projective K3 surface (so $K_S \cong \mathcal{O}_S$,
$h^{1,0}(S) = 0$, $h^{2,0}(S) = 1$; \cite{Mukai1984}),
$E$ is a smooth elliptic curve over $\C$ with marked origin
$0_E \in E$, and
\[
X := S \times E.
\]
The canonical bundle is trivial, $K_X \cong p_S^*K_S \otimes p_E^*K_E
\cong \mathcal O_X$, and Hodge invariants are
\[
\dim_\C X = 3,\qquad
h^{1,0}(X)=h^{1,0}(E)=1,\qquad
\chi_{\mathrm{top}}(X) = \chi_{\mathrm{top}}(S)\cdot\chi_{\mathrm{top}}(E)
= 24 \cdot 0 = 0,\qquad
\chi(\mathcal O_X) = \chi(\mathcal O_S)\cdot\chi(\mathcal O_E) = 2\cdot 0 = 0.
\]
The threefold $X$ is Calabi--Yau in the trivial-canonical sense used by
Donaldson--Thomas theory and shifted symplectic geometry; the
construction uses the chosen trivialization of \(K_X\).
\end{definition}

\begin{remark}[Distinguishing $\chi_{\mathrm{top}}(K3) = 24$ and
$\chi(\mathcal O_S) = 2$]
\label{rem:two-K3-Eulers}
The two integers enter the construction by different routes.  The
Borcherds weight is $f(0,0)/2 = 5$, where $f(0,0) = 10$ is the
constant Fourier coefficient of the weight-zero index-one weak Jacobi
form $\phi_{0,1} = r^{-1} + 10 + r + O(q)$ in the Eichler--Zagier
normalization (computed, \cite{EZ:1}).  The topological Euler
number $\chi_{\mathrm{top}}(K3) = 24$ enters one step removed, through
the elliptic-genus specialization: at $z = 0$ the $\phi_{0,1}$ support
$4n - l^2 \ge -1$ leaves only $l \in \{-1,0,1\}$ at $n = 0$, so the
constant term is $\sum_l f(0,l) = f(0,-1) + f(0,0) + f(0,1) = 12$, and
$Z_{K3}(\tau,0) = 2\,\phi_{0,1}(\tau,0)$ has constant term
$2 \cdot 12 = 24 = \chi_{\mathrm{top}}(K3)$.  Thus $\chi_{\mathrm{top}}(K3)
= 2\sum_l f(0,l)$, not $f(0,0)$; the Borcherds weight $5$ sees only the
single coefficient $f(0,0) = 10$.  The integer
$\chi(\mathcal O_S) = 2$ enters by a third route, the categorical
Calabi--Yau dimension and the holomorphic part of the Mukai pairing.
The two integers belong to different entries: the $\kappa$-tuple of \cite{CYQGVolIII}
records $\kappa_{\mathrm{fiber}} = 24 = \chi_{\mathrm{top}}(K3)$, not
$2 = \chi(\mathcal O_S)$, and the universal Borcherds identity
$\kappa_{\mathrm{BKM}}(\Phi_N) = c_N(0)/2$ at $\Phi_1 = \Delta_5$ gives
$c_1(0) = 10$ and $\kappa_{\mathrm{BKM}} = 5$.  The expression
$\kappa_{\mathrm{ch}}^{\mathrm{Hodge}}+\chi(\mathcal O_{K3})=0+2=2$
does not give the Borcherds weight, while
$\kappa_{\mathrm{ch}}^{\mathrm{Heis}}+\chi(\mathcal O_{K3})=3+2=5$
is only an accidental equality at the \(N=1\) row.  Across the CHL
frame the Borcherds weights are \(c_N(0)/2=(5,3,2,\tfrac32,1)\)
\textup{(Jatkar--Sen \cite{JS06Dyon}; Govindarajan--Krishna
\cite{GK09CHL,GK10Composite}; the \(N=4\) weight is
half-integral)}, not a fixed
sum of a chiral rank and \(\chi(\mathcal O_{K3})\)
(proved, \cite{CYQGVolIII} counterexample).  The coordinate $\kappa_{\mathrm{BKM}} = 5$ is the
cusp-form weight; the protected scalar shadow is written at level
$n = \chi(\mathcal O_Y)$ for the one-dimensional subscheme
$Y\subset K3\times E$.  The target $X = K3\times E$ is a
trivial-canonical threefold.
\end{remark}

\subsection{Mukai lattice on the K3 factor}
\label{subsec:k3-mukai-lattice}

\begin{definition}[Mukai lattice and Mukai pairing]
\label{def:mukai-lattice}
Let $S$ be a smooth complex K3 surface.  The \emph{Mukai lattice} is the
total even cohomology
\[
\Lambda_S
:= \widetilde H(S, \Z)
= H^0(S, \Z) \oplus H^2(S, \Z) \oplus H^4(S, \Z),
\]
identified with the K-theoretic charge lattice
$K_0(D^b(\mathrm{Coh}\,S))_{\mathrm{tor.\,free}}$ via the Mukai vector
\[
v(\mathcal F)
:= \mathrm{ch}(\mathcal F)\sqrt{\mathrm{Td}(S)}
= \bigl(r(\mathcal F),\,c_1(\mathcal F),\,r(\mathcal F) + \mathrm{ch}_2(\mathcal F)\bigr),
\]
where $r$ is the rank, $c_1$ is the first Chern class, and the half-shift
$r$ in the four-form term comes from $\mathrm{Td}(S) = 1 + 2[\mathrm{pt}]$
(proved, \cite{Mukai1984,Mukai1987}).  The \emph{Mukai pairing}
is
\[
\bigl\langle (r, c, s),\,(r', c', s')\bigr\rangle_{\mathrm{Muk}}
:= c\cdot c' - r s' - r' s,
\]
where $c\cdot c' \in \Z$ is the cup product on $H^2(S, \Z)$.  It is even,
unimodular, of signature $(4, 20)$, and rank $24$:
\[
\Lambda_S \cong II_{4, 20}
\cong U^{\oplus 4} \oplus E_8(-1)^{\oplus 2},
\]
where $U$ denotes the rank-two hyperbolic plane $U \cong \mathbb Z e
\oplus \mathbb Z f$ with $e^2 = f^2 = 0$, $e\cdot f = 1$, and $E_8(-1)$
denotes the negative-definite $E_8$ root lattice (proved,
\cite{Mukai1987,NIKULIN}).
\end{definition}

\begin{remark}[Sign convention]
\label{rem:mukai-sign-convention}
The signature $(4, 20)$ uses the convention that $H^0\oplus H^4$
contributes a hyperbolic plane $U$ via $(1, 0, 0)$ and $(0, 0, 1)$ and
that $H^2(S, \Z) \cong II_{3, 19} \cong U^{\oplus 3} \oplus
E_8(-1)^{\oplus 2}$, which has positive part of dimension $3$.  Together,
$\Lambda_S = U \oplus II_{3, 19} \cong U^{\oplus 4} \oplus
E_8(-1)^{\oplus 2}$ with positive part of dimension $4$ (proved,
\cite{Mukai1987}).  The signature is $(4, 20)$, not $(20, 4)$: the
positive direction is the smaller side.  A
reflection $r_v$ across $v\in\Lambda_S$ acts by $r_v(x) = x - 2(x, v)v/(v,
v)$, taking the $\R$-linear extension when $(v, v)\ne 0$, without
implicit sign reversal.
\end{remark}

\subsection{K\"unneth and the formal even Mukai sector on $X = S \times E$}
\label{subsec:kunneth-formal-mukai}

\begin{definition}[Formal even Mukai sector on $X = S \times E$]
\label{def:formal-even-mukai-sector}
Pick a basepoint $0_E \in E$ and the standard generators
$1_E \in H^0(E, \Z)$ and $\omega_E \in H^2(E, \Z)$ with $\int_E \omega_E
= 1$.  Every even cohomology class on $X = S \times E$ decomposes as
\[
v_X = P \otimes 1_E + Q \otimes \omega_E,\qquad
P, Q \in \Lambda_S
\]
(proved, K\"unneth on simply connected $\widetilde H$ tensor
$E$-cohomology).  The pair $(Q, P)\in\Lambda_S\oplus\Lambda_S$ is the
\emph{formal even charge}.  The component $P\otimes 1_E$ is invariant
under translation along $E$, while $Q\otimes\omega_E$ records the
$E$-degree.

The \emph{formal full-Mukai dyonic lattice} is the additive abelian
group
\[
\Gamma_X^{\mathrm{form}} := \Lambda_S \oplus \Lambda_S,\qquad
c = (Q, P),
\]
with electric component $Q$ and magnetic component $P$.  The symbol
\(\Gamma_X^{\mathrm{phys}}\) is defined here by the equality
\[
\Gamma_X^{\mathrm{phys}} := \Gamma_X^{\mathrm{form}}.
\]
This is not the full four-dimensional $\mathcal N = 4$ electric--magnetic
charge lattice and not compact Hall support.  It is the D6--D2--D0/OP
even branch on which the Mukai Gram map is tested before imposing
algebraicity, effectivity, stability, or non-emptiness of compact Hall
moduli.
\end{definition}

\begin{definition}[Liu product-stability datum on \(S\times E\)]
\label{def:liu-product-stability-datum}
Fix a rational Bridgeland stability condition
\[
\sigma_S=(\mathcal A_S,Z_S),\qquad
Z_S\bigl(K(D^b(\mathrm{Coh}\,S))\bigr)\subset\Q+i\Q ,
\]
on \(D^b(\mathrm{Coh}\,S)\).  Let \(L_E\) be a degree-one line bundle
on the elliptic curve \(E\), and let
\[
p:S\times E\to S,\qquad q:S\times E\to E
\]
be the projections.  The Abramovich--Polishchuk heart used by Liu is
the following heart
\cite{AbramovichPolishchukTStructures,LiuProductStability}:
\[
\mathcal A_E^{\mathrm{AP}}
=\bigl\{M\in D^b(\mathrm{Coh}(S\times E))\mid
p_*(M\otimes q^*L_E^{\otimes n})\in\mathcal A_S
\text{ for }n\gg0\bigr\}.
\]
For \(M\in\mathcal A_E^{\mathrm{AP}}\), Liu's curve-factor
polynomial is
\[
L_M(n):=
Z_S\!\bigl(p_*(M\otimes q^*L_E^{\otimes n})\bigr)
=a(M)n+b(M)+i\bigl(c(M)n+d(M)\bigr),
\]
where \(a,b,c,d:K(\mathcal A_E^{\mathrm{AP}})\to\Q\) are additive
homomorphisms.  For \(t\in\Q_{>0}\), set
\[
Z_t(M)=a(M)t-d(M)+i\,c(M)t,\qquad
\nu_t(M)=
\begin{cases}
\dfrac{-a(M)t+d(M)}{c(M)t},&c(M)\neq0,\\[6pt]
+\infty,&c(M)=0.
\end{cases}
\]
Let
\[
\mathcal T_t=\{M\in\mathcal A_E^{\mathrm{AP}}\mid
\nu_{t,\min}(M)>0\},\qquad
\mathcal F_t=\{M\in\mathcal A_E^{\mathrm{AP}}\mid
\nu_{t,\max}(M)\le0\},
\]
and
\[
\mathcal A_E^t=\langle\mathcal T_t,\mathcal F_t[1]\rangle .
\]
For \(s\in\Q_{>0}\), Liu's product central charge is
\[
Z_E^{s,t}(M)=c(M)s+b(M)+i\bigl(-a(M)t+d(M)\bigr).
\]
The Liu product-stability condition is
\[
\sigma_{s,t}:=(\mathcal A_E^t,Z_E^{s,t}).
\]
Liu proves that this is a rational Bridgeland stability condition for
\(s,t\in\Q_{>0}\), and extends the construction continuously to
\((s,t)\in\R_{>0}^2\)
\cite[Theorem~4.7 and \S5]{LiuProductStability}.  This is a
source-imported stability condition.  It is not a finite-type moduli
construction, not a compact Hall stage, and not a Pfaffian
orientation.
\end{definition}

The packet \(\texttt{certificates/moduli/liu\_product\_stability}\)
checks this definition.  It records the source theorem, the AP heart
membership formula, the linear polynomial \(L_M(n)\), the four
coefficient homomorphisms \(a,b,c,d\), the weak charge \(Z_t\), the
torsion-pair tilt, and the central charge \(Z_E^{s,t}\).  Its verifier
status is \texttt{LIU\_PRODUCT\_STABILITY\_DEFINED}.  It does not
prove the separate AP-compatibility row, finite-type semistable
substacks, derived enhancements, compact Hall stages, Pfaffian
orientations, or protected traces.

\begin{proposition}[Compatibility with the Abramovich--Polishchuk heart]
\label{prop:ap-liu-heart-compatibility}
\textnormal{[proved]}
With the datum of
Definition~\ref{def:liu-product-stability-datum}, the heart of Liu's
product stability condition is the torsion-pair tilt of the
Abramovich--Polishchuk heart:
\[
\heartsuit(\sigma_{s,t})
=\mathcal A_E^t
=\langle\mathcal T_t,\mathcal F_t[1]\rangle,
\]
where \((\mathcal T_t,\mathcal F_t)\) is the torsion pair in
\(\mathcal A_E^{\mathrm{AP}}\) defined by the slope \(\nu_t\).
Consequently
\[
\mathcal A_X^{\mathrm{AP/Liu}}=\mathcal A_E^t.
\]
Thus the retained heart is Abramovich--Polishchuk-compatible in the
finite-row sense: it is obtained from \(\mathcal A_E^{\mathrm{AP}}\)
by Liu's specified torsion-pair tilt.
\end{proposition}

\begin{proof}
The Abramovich--Polishchuk heart
\(\mathcal A_E^{\mathrm{AP}}\) is the global heart entering Liu's
construction.  Liu defines the weak charge \(Z_t\) on this heart,
uses the resulting slope \(\nu_t\) to form the torsion pair
\[
\mathcal T_t=\{M\mid\nu_{t,\min}(M)>0\},\qquad
\mathcal F_t=\{M\mid\nu_{t,\max}(M)\le0\},
\]
and then sets
\[
\mathcal A_E^t=\langle\mathcal T_t,\mathcal F_t[1]\rangle
\]
before defining
\[
\sigma_{s,t}=(\mathcal A_E^t,Z_E^{s,t})
\]
\cite[\S2.2 and Theorem~4.7]{LiuProductStability}.  Hence the heart
of \(\sigma_{s,t}\) is \(\mathcal A_E^t\) by construction.  The
retained heart is defined as
\(\mathcal A_X^{\mathrm{AP/Liu}}:=\heartsuit(\sigma_{s,t})\), so it
equals \(\mathcal A_E^t\).  This proves the compatibility statement
using only the definition of the tilt; it proves no noetherianity,
HN-filtration, boundedness, finite-moduli, derived-enhancement,
Pfaffian-orientation, or protected-trace row.
\end{proof}

The packet
\(\texttt{certificates/moduli/ap\_liu\_compatibility}\) checks this
heart identification.  It records the AP membership formula, the
weak charge \(Z_t\), the slope \(\nu_t\), the torsion pair
\((\mathcal T_t,\mathcal F_t)\), the tilted heart
\(\mathcal A_E^t\), and the equality
\(\mathcal A_X^{\mathrm{AP/Liu}}=\heartsuit(\sigma_{s,t})=\mathcal A_E^t\).
Its verifier status is
\texttt{AP\_LIU\_COMPATIBILITY\_VERIFIED}.  It does not prove
noetherianity, HN filtrations, bounded HN types, finite-type
semistable substacks, derived enhancements, compact Hall stages,
Pfaffian orientations, or protected traces.

\begin{definition}[Retained Abramovich--Polishchuk/Liu heart]
\label{def:retained-ap-liu-heart}
Fix the Liu product-stability datum of
Definition~\ref{def:liu-product-stability-datum}.  The retained
stability heart on
\[
D^b(\mathrm{Coh}\,X),\qquad X=S\times E,
\]
is
\[
\mathcal A_X^{\mathrm{AP/Liu}}
:=\heartsuit(\sigma_{s,t}),
\]
equivalently
\(\mathcal A_X^{\mathrm{AP/Liu}}=\mathcal A_E^t\) by
Proposition~\ref{prop:ap-liu-heart-compatibility}.  In
the rest of the manuscript the symbol \(\sigma\) means this retained
datum \(\sigma_{s,t}\), and its heart is
\(\mathcal A_X^{\mathrm{AP/Liu}}\).

This definition is only the choice of the ambient heart and stability
datum.  Noetherianity of the retained windows, existence of
Harder--Narasimhan filtrations, boundedness of HN types, finite-type
semistable substacks, quasi-smooth derived enhancements, universal
complexes, and compact Hall stages are separate finite rows.
\end{definition}

The packet \(\texttt{certificates/moduli/stability\_heart}\) checks
this definition as a table-level datum.  It records the inputs
\[
(S,H),\quad (E,0_E),\quad \sigma_S,\quad L_E,\quad
\mathcal A_X^{\mathrm{AP}},\quad \sigma_{s,t},\quad
\mathcal R_{\mathrm{HN}},
\]
and verifies the scalar firewall excluding noetherianity, HN
filtrations, bounded HN types, finite-type moduli, Pfaffian
orientations, and protected traces.  Its verifier status is
\texttt{STABILITY\_HEART\_DATUM\_VERIFIED}.

\begin{definition}[Retained noetherian window]
\label{def:retained-noetherian-window}
Fix the heart \(\mathcal A_X^{\mathrm{AP/Liu}}\).  A retained finite
window at height \(R\) is a full exact subcategory
\[
\mathcal A_{X,R}^{\mathrm{ret}}\subset\mathcal A_X^{\mathrm{AP/Liu}}
\]
with a finite set of isomorphism classes, closed under retained
subobjects and retained quotients, together with a length function
\[
\ell_R:\operatorname{Obj}(\mathcal A_{X,R}^{\mathrm{ret}})\to\Z_{\ge0}
\]
such that \(\ell_R(A)=0\) exactly for the zero object and
\(\ell_R(A)<\ell_R(B)\) for every strict retained monomorphism
\(A\hookrightarrow B\).
\end{definition}

\begin{proposition}[Retained finite windows are noetherian]
\label{prop:retained-window-noetherian}
\textnormal{[proved]}
Every retained finite window
\(\mathcal A_{X,R}^{\mathrm{ret}}\) in the sense of
Definition~\ref{def:retained-noetherian-window} is noetherian: every
ascending chain of retained subobjects terminates.
\end{proposition}

\begin{proof}
Let
\[
A_0\subset A_1\subset A_2\subset\cdots
\]
be a chain of retained subobjects of an object \(A\) in
\(\mathcal A_{X,R}^{\mathrm{ret}}\).  After deleting repetitions, each
inclusion is strict.  The length function then gives a strictly
increasing sequence
\[
\ell_R(A_0)<\ell_R(A_1)<\ell_R(A_2)<\cdots .
\]
All \(A_i\) lie in the finite retained window and all are subobjects of
\(A\), so \(\ell_R(A_i)\le \ell_R(A)\).  A strictly increasing sequence
of non-negative integers bounded above by \(\ell_R(A)\) is finite.
Thus the original ascending chain stabilizes.
\end{proof}

The packet
\(\texttt{certificates/moduli/retained\_window\_noetherian}\) checks
this finite ACC statement on a displayed retained exact window.  It
verifies a finite object set, subobject and quotient closure, strict
increase of the length rank along each strict monomorphism, acyclicity
of the strict subobject graph, and the longest-chain bound.  Its
verifier status is
\texttt{RETAINED\_WINDOW\_NOETHERIAN\_VERIFIED}.  It does not prove
global noetherianity of \(\mathcal A_X^{\mathrm{AP/Liu}}\),
Harder--Narasimhan filtrations, bounded HN types, finite-type
semistable substacks, quasi-smooth moduli, Pfaffian orientations, or
protected traces.

\begin{definition}[Retained finite Harder--Narasimhan filtration datum]
\label{def:retained-finite-hn-filtration-datum}
Let \(\mathcal A_{X,R}^{\mathrm{ret}}\) be a retained finite window.
A retained finite Harder--Narasimhan filtration datum consists of:
\[
\Phi_R\subset\Q,\qquad
\phi_R:\mathcal S_R\to\Phi_R,
\]
where \(\Phi_R\) is a finite totally ordered phase set and
\(\mathcal S_R\) is a class of nonzero retained semistable objects,
together with, for every \(A\in\mathcal A_{X,R}^{\mathrm{ret}}\), a
finite exact filtration
\[
0=A_0\subset A_1\subset\cdots\subset A_m=A
\]
whose retained quotients \(E_i=A_i/A_{i-1}\) lie in \(\mathcal S_R\)
and satisfy
\[
\phi_R(E_1)>\phi_R(E_2)>\cdots>\phi_R(E_m).
\]
The zero object has the empty filtration.  This is finite
retained-window data.  It is not a global Harder--Narasimhan theorem
for \(\mathcal A_X^{\mathrm{AP/Liu}}\), not boundedness of HN types,
and not finite-type moduli.
\end{definition}

\begin{proposition}[Retained finite HN filtrations exist from the datum]
\label{prop:retained-window-hn-filtration}
\textnormal{[proved]}
If a retained finite window
\(\mathcal A_{X,R}^{\mathrm{ret}}\) carries the datum of
Definition~\ref{def:retained-finite-hn-filtration-datum}, then every
object of \(\mathcal A_{X,R}^{\mathrm{ret}}\) has a retained
Harder--Narasimhan filtration.
\end{proposition}

\begin{proof}
For \(A=0\) the empty filtration is the Harder--Narasimhan filtration.
For \(A\neq0\), the datum gives an exact retained chain
\[
0=A_0\subset A_1\subset\cdots\subset A_m=A .
\]
Each quotient \(E_i=A_i/A_{i-1}\) is retained and semistable, and the
phase inequalities are strict.  The chain is finite, hence terminates
inside the exact window.  These are precisely the
Harder--Narasimhan conditions in the retained finite window.
\end{proof}

The packet
\(\texttt{certificates/moduli/retained\_window\_hn\_filtration}\)
checks this finite statement on a displayed retained exact window.  It
verifies a finite phase set, semistable quotient factors, exact
filtration chains, length additivity, and strict phase descent.  Its
verifier status is
\texttt{RETAINED\_WINDOW\_HN\_FILTRATION\_VERIFIED}.  It does not
prove global Harder--Narasimhan filtrations in
\(\mathcal A_X^{\mathrm{AP/Liu}}\), bounded HN types, finite-type
semistable substacks, quasi-smooth moduli, Pfaffian orientations, or
protected traces.

\begin{definition}[Retained bounded HN type datum]
\label{def:retained-bounded-hn-type-datum}
Let \(\mathcal A_{X,R}^{\mathrm{ret}}\) be a retained finite window
with a retained finite Harder--Narasimhan filtration datum.  The
\emph{retained HN type} of an object \(A\) is the finite word
\[
\operatorname{HNType}_R(A)
=\bigl(([E_1],\phi_R(E_1)),\ldots,([E_m],\phi_R(E_m))\bigr)
\]
attached to its retained HN quotients \(E_i=A_i/A_{i-1}\), with the
empty word for \(A=0\).  A retained bounded HN type datum is a finite
set \(\mathfrak T_R\) containing all these words, together with a
finite Liu--Hilbert decoration of the semistable factor classes:
\[
[E]\longmapsto
\bigl(\widehat c_E,[a_E,b_E],P_E,N_E,\mathcal F_E\bigr).
\]
Here \(\widehat c_E\) is the retained numerical class, \([a_E,b_E]\)
is the cohomological-amplitude interval, \(P_E\) is the retained
Hilbert-polynomial label, \(N_E\) is the retained regularity bound, and
\(\mathcal F_E\) is the retained closure family.  This datum is a
finite bound on active-window HN types.  It is not a finite-type
semistable substack, not a quasi-smooth derived enhancement, and not a
compact Hall stage.
\end{definition}

\begin{proposition}[HN type is bounded in a retained active window]
\label{prop:retained-hn-type-bounded}
\textnormal{[proved]}
Let \(\mathcal A_{X,R}^{\mathrm{ret}}\) be a retained finite window
with a retained finite Harder--Narasimhan filtration datum.  Then the
set of retained HN types occurring in
\(\mathcal A_{X,R}^{\mathrm{ret}}\) is finite.  More precisely, if
\[
N_R=\max\{\ell_R(A)\mid A\in
\mathcal A_{X,R}^{\mathrm{ret}}\},
\]
then every retained HN type has length at most \(N_R\).
\end{proposition}

\begin{proof}
The retained window has finitely many isomorphism classes, hence the
set of retained semistable HN factors is finite.  If
\[
0=A_0\subset A_1\subset\cdots\subset A_m=A
\]
is the retained HN filtration of \(A\), then each inclusion
\(A_{i-1}\subset A_i\) is strict unless \(m=0\).  The retained length
rank gives
\[
0=\ell_R(A_0)<\ell_R(A_1)<\cdots<\ell_R(A_m)\le N_R,
\]
so \(m\le N_R\).  Thus every retained HN type is a word of length at
most \(N_R\) in a finite alphabet of semistable factor classes.  The
set of such words is finite, and the strict phase condition selects a
finite subset.  A finite Liu--Hilbert decoration of the factor classes
therefore gives a finite decorated HN type set.
\end{proof}

The packet \(\texttt{certificates/moduli/retained\_hn\_type\_bounds}\)
checks this finite boundedness statement.  It verifies the finite
semistable factor-type set, the finite HN type-word set, the maximal
HN height, the maximal factor count, the retained length-rank bound,
and the mirror of the HN type rows in
\(\texttt{certificates/moduli/k3e\_finite\_moduli/hn\_type\_bounds.csv}\).
Its verifier status is
\texttt{RETAINED\_HN\_TYPE\_BOUNDS\_VERIFIED}.  It does not construct
finite-type semistable substacks, quasi-smooth derived enhancements,
universal complexes, compact Hall stages, Pfaffian orientations, or
protected traces.

\begin{definition}[Retained finite-type semistable substack datum]
\label{def:retained-finite-type-semistable-substack-datum}
Fix a retained finite window
\(\mathcal A_{X,R}^{\mathrm{ret}}\), the retained stability condition
\(\sigma_{s,t}\), and the retained bounded HN type datum of
Definition~\ref{def:retained-bounded-hn-type-datum}.  A retained
finite-type semistable substack datum assigns to every retained
semistable factor type \(c\) a finite-type bounded
Quot/Postnikov ambient stack \(Q_{R,c}\), a chart
\[
U_{R,c}\to Q_{R,c},
\]
and a specialization-closed finite-type algebraic substack
\[
\mathfrak M^{ss}_{R,c}\subset Q_{R,c}
\]
whose objects are exactly the \(\sigma_{s,t}\)-semistable retained
objects of type \(c\).  The datum carries two defect ranks:
\[
d^{ss}_{R,c}=0,\qquad d^{bd}_{R,c}=0.
\]
The first says that no unstable object is included and no semistable
object of type \(c\) is omitted.  The second says that the bounded
Quot/Postnikov ambient covers the retained class.  This is a
finite-type semistable-substack datum.  It is not a quasi-smooth
derived enhancement, not a universal complex, not a rigidification,
not a compact Hall correspondence, and not an orientation datum.
\end{definition}

\begin{proposition}[Finite-type semistable substacks for retained classes]
\label{prop:retained-finite-type-semistable-substacks}
\textnormal{[proved]}
Let \(\mathcal A_{X,R}^{\mathrm{ret}}\) carry a retained bounded HN
type datum and a retained finite-type semistable substack datum.  Then
for every retained semistable factor type \(c\) the stack
\(\mathfrak M^{ss}_{R,c}\) is of finite type over \(\C\), is
specialization-closed inside its bounded Quot/Postnikov ambient, and
parametrizes exactly the retained \(\sigma_{s,t}\)-semistable objects
of type \(c\).
\end{proposition}

\begin{proof}
The retained HN type datum gives only finitely many semistable factor
types \(c\).  For each such \(c\), the semistable-substack datum places
\(\mathfrak M^{ss}_{R,c}\) inside a finite-type bounded
Quot/Postnikov ambient \(Q_{R,c}\) through the chart \(U_{R,c}\), with
the semistable locus represented by finitely many locally closed
retained chart pieces.  Finite unions of locally closed substacks of a
finite-type algebraic stack are finite type.  The condition
\(d^{bd}_{R,c}=0\) says that the ambient covers the retained class,
while \(d^{ss}_{R,c}=0\) says that the substack is precisely the
\(\sigma_{s,t}\)-semistable locus of that class.  The
specialization-closed flag is part of the datum and is checked in the
same finite chart.  Thus each \(\mathfrak M^{ss}_{R,c}\) is finite
type and has the asserted moduli meaning.  The argument uses no
derived enhancement, universal complex, rigidification, finite-inertia
stratification, flag correspondence, cosection atlas, transition map,
Pfaffian orientation, or protected trace.
\end{proof}

The packet
\(\texttt{certificates/moduli/retained\_semistable\_substacks}\)
checks this finite-type row.  It imports the retained HN factor types,
supplies one bounded Quot/Postnikov ambient chart and one semistable
substack for each retained class, verifies specialization-closedness,
and checks zero semistability and boundedness defect ranks.  Its
verifier status is
\texttt{RETAINED\_SEMISTABLE\_SUBSTACKS\_VERIFIED}.  The mirrored
rows in
\(\texttt{certificates/moduli/k3e\_finite\_moduli/semistable\_substacks.csv}\)
discharge item \(167\).  They do not construct derived enhancements,
universal complexes, scalar rigidifications, finite inertia strata,
extension or flag stacks, cosection atlases, transitions, compact Hall
stages, Pfaffian orientations, or protected traces.

\begin{definition}[Retained quasi-smooth derived enhancement datum]
\label{def:retained-quasi-smooth-derived-enhancement-datum}
Fix the retained finite-type semistable substacks
\(\mathfrak M^{ss}_{R,c}\) of
Definition~\ref{def:retained-finite-type-semistable-substack-datum}.
A retained quasi-smooth derived enhancement datum assigns to every
retained semistable factor type \(c\) a derived Artin stack
\[
\mathfrak M^{\mathrm{der}}_{R,c}
\]
with a morphism
\(\mathfrak M^{\mathrm{der}}_{R,c}\to\operatorname{Perf}(X)\), an
identification
\[
t_0\mathfrak M^{\mathrm{der}}_{R,c}
=\mathfrak M^{ss}_{R,c},
\]
a perfect cotangent complex of amplitude
\[
\mathbb L_{\mathfrak M^{\mathrm{der}}_{R,c}}\in
D^{[-1,0]}(\mathfrak M^{\mathrm{der}}_{R,c}),
\]
and a restricted PTVV \((-1)\)-shifted symplectic form
\[
\omega^{\mathrm{PTVV}}_{R,c}.
\]
The datum carries two defect ranks
\[
d^{\mathrm{Tor}}_{R,c}=0,\qquad d^\omega_{R,c}=0.
\]
The first says that the cotangent-amplitude assertion has no retained
Tor-amplitude defect.  The second says that the PTVV form restricts to
the retained derived substack without symplectic defect.  This is a
derived-enhancement datum.  It is not a universal perfect complex, not
a scalar rigidification, not an inertia stratification, not a
Hall-flag correspondence, and not a cosection or orientation datum.
\end{definition}

\begin{proposition}[Quasi-smooth derived Artin enhancements]
\label{prop:retained-quasi-smooth-derived-enhancements}
\textnormal{[proved]}
Let the retained finite-type semistable substacks carry the datum of
Definition~\ref{def:retained-quasi-smooth-derived-enhancement-datum}.
Then every \(\mathfrak M^{ss}_{R,c}\) has a quasi-smooth derived Artin
enhancement \(\mathfrak M^{\mathrm{der}}_{R,c}\) whose classical
truncation is \(\mathfrak M^{ss}_{R,c}\), whose cotangent complex has
amplitude \([-1,0]\), and whose \((-1)\)-shifted symplectic form is
the restricted PTVV form on the Calabi--Yau threefold
\(X=S\times E\).
\end{proposition}

\begin{proof}
The datum gives a derived Artin stack
\(\mathfrak M^{\mathrm{der}}_{R,c}\) with
\(t_0\mathfrak M^{\mathrm{der}}_{R,c}=\mathfrak M^{ss}_{R,c}\).
The finite-type statement for the classical truncation is
Proposition~\ref{prop:retained-finite-type-semistable-substacks}.  The
cotangent-amplitude condition
\(\mathbb L_{\mathfrak M^{\mathrm{der}}_{R,c}}\in D^{[-1,0]}\) and the
vanishing \(d^{\mathrm{Tor}}_{R,c}=0\) are exactly the
quasi-smoothness criterion for a derived Artin stack.  Since
\(K_X\simeq\mathcal O_X\), \(X\) is a Calabi--Yau threefold, and PTVV
produce the ambient \((-1)\)-shifted symplectic form on the derived
moduli of perfect complexes \cite[Theorem~0.3]{PTVV2013}.  The
vanishing \(d^\omega_{R,c}=0\) says that this form restricts to the
retained derived substack.  Thus the retained enhancement is
quasi-smooth and carries the asserted PTVV form.  The proof supplies no
universal complex on the rigidified stack, no scalar rigidification, no
finite-inertia stratification, no extension or flag stack, no cosection
atlas, no transition system, no Pfaffian orientation, and no protected
trace.
\end{proof}

The packet
\(\texttt{certificates/moduli/retained\_derived\_enhancements}\)
checks this row.  It imports the finite-type semistable substacks,
records one derived Artin enhancement per retained substack, verifies
cotangent amplitude \([-1,0]\), imports the PTVV \((-1)\)-shifted
symplectic form, and checks zero Tor-amplitude and symplectic defect
ranks.  Its verifier status is
\texttt{RETAINED\_DERIVED\_ENHANCEMENTS\_VERIFIED}.  The mirrored rows
in
\(\texttt{certificates/moduli/k3e\_finite\_moduli/derived\_enhancements.csv}\)
discharge item \(168\).  They do not construct universal complexes,
scalar rigidifications, finite inertia strata, extension or flag
stacks, cosection atlases, transitions, compact Hall stages, Pfaffian
orientations, or protected traces.

\begin{definition}[Retained scalar rigidification and residual-inertia datum]
\label{def:retained-rigidification-inertia-datum}
Fix the retained quasi-smooth derived enhancements
\(\mathfrak M^{\mathrm{der}}_{R,c}\) of
Definition~\ref{def:retained-quasi-smooth-derived-enhancement-datum}.
A retained scalar rigidification and residual-inertia datum assigns to
every retained semistable factor type \(c\) an exact sequence of
stabilizer groups
\[
1\longrightarrow \mathbb G_m\longrightarrow
\operatorname{Aut}_{R,c}\longrightarrow H_{R,c}\longrightarrow 1,
\]
a scalar rigidification
\[
\mathfrak M^{\mathrm{rig}}_{R,c}
=\bigl[\mathfrak M^{\mathrm{der}}_{R,c}/\mathbb G_m\bigr],
\]
and the residual inertia group \(H_{R,c}\).  The datum requires
\[
|H_{R,c}|<\infty,\qquad
d^{\mathrm{rig}}_{R,c}=0,\qquad d^{\mathrm{in}}_{R,c}=0.
\]
Thus the scalar \(\mathbb G_m\) has been removed, the residual inertia
is finite, and the rigidification and inertia defects vanish.  This is
a scalar-rigidification datum for retained derived substacks.  It is
not a universal perfect complex, not a finite closed inertia
stratification, not an extension or flag stack, not a cosection atlas,
not a transition system, and not an orientation datum.
\end{definition}

\begin{proposition}[Finite residual inertia after scalar rigidification]
\label{prop:retained-finite-residual-inertia-after-rigidification}
\textnormal{[proved]}
Let the retained derived enhancements carry the datum of
Definition~\ref{def:retained-rigidification-inertia-datum}.  Then for
each retained semistable factor type \(c\), the scalar rigidified stack
\(\mathfrak M^{\mathrm{rig}}_{R,c}\) has finite residual inertia
\(H_{R,c}\).  Equivalently, after quotienting the scalar
\(\mathbb G_m\)-automorphisms, no positive-dimensional stabilizer
remains in the retained row.
\end{proposition}

\begin{proof}
The datum gives the exact sequence
\[
1\to \mathbb G_m\to\operatorname{Aut}_{R,c}\to H_{R,c}\to 1.
\]
The scalar rigidification is the quotient by the indicated central
\(\mathbb G_m\).  Therefore the stabilizer remaining on
\(\mathfrak M^{\mathrm{rig}}_{R,c}\) is \(H_{R,c}\).  The condition
\(|H_{R,c}|<\infty\) gives finite residual inertia, and the equalities
\(d^{\mathrm{rig}}_{R,c}=d^{\mathrm{in}}_{R,c}=0\) say that neither
the quotient nor the inertia computation has a retained defect.  The
argument constructs no universal complex, finite closed inertia
stratification, flag stack, cosection atlas, transition map, Pfaffian
orientation, or protected trace.
\end{proof}

The packet
\(\texttt{certificates/moduli/retained\_rigidification\_inertia}\)
checks this row.  It imports the retained derived enhancements,
records one scalar rigidification for each retained substack, verifies
the exact sequence
\[
1\to\mathbb G_m\to\operatorname{Aut}_{R,c}\to H_{R,c}\to 1,
\]
checks \(\mathbb G_m\)-removal, finite residual inertia, and zero
rigidification and inertia defect ranks.  Its verifier status is
\texttt{RETAINED\_RIGIDIFICATION\_INERTIA\_VERIFIED}.  The mirrored
rows in
\(\texttt{certificates/moduli/k3e\_finite\_moduli/scalar\_rigidifications.csv}\)
discharge item \(169\) and the scalar-rigidification existence row
item \(175\).  They do not construct universal complexes, finite
closed inertia stratifications, extension or flag stacks, cosection
atlases, transitions, compact Hall stages, Pfaffian orientations, or
protected traces.

\begin{definition}[Retained finite class-bound datum]
\label{def:retained-finite-class-bound-datum}
Fix the retained finite-type semistable substacks
\(\mathfrak M^{ss}_{R,c}\) and the retained HN type-bound rows.  A
retained finite class-bound datum consists of a finite set
\[
C_R=\{c_1,\ldots,c_m\}
\]
of retained semistable classes, together with maps assigning to each
\(c\in C_R\) a retained type \(\tau(c)\), a charge \(\widehat c\), a
substack \(\mathfrak M^{ss}_{R,c}\), a finite Hilbert-polynomial label
set
\[
\{P_{c,i}\}_{i\in I_c},\qquad |I_c|<\infty,
\]
a cohomological-amplitude interval
\[
[a_c,b_c]\subset\mathbb Z,\qquad a_c\le b_c,
\]
and a nonnegative regularity bound \(N_c\).  The datum requires zero
class-coverage, Hilbert-bound, amplitude-bound, and regularity-bound
defects.  This is a retained class-bound datum.  It is not a universal
perfect complex, not a finite closed inertia stratification, not an
extension or flag stack, not a cosection atlas, not a transition
system, and not an orientation datum.
\end{definition}

\begin{proposition}[Finite retained class bounds]
\label{prop:retained-finite-class-bounds}
\textnormal{[proved]}
Let a retained active window carry the datum of
Definition~\ref{def:retained-finite-class-bound-datum}.  Then the
retained class set \(C_R\) is finite.  For every \(c\in C_R\), the
Hilbert-polynomial labels \(P_{c,i}\), the cohomological amplitude
interval \([a_c,b_c]\), and the regularity bound \(N_c\) are finite
retained data.
\end{proposition}

\begin{proof}
The datum writes \(C_R=\{c_1,\ldots,c_m\}\), hence the retained class
set is finite.  For each \(c\), the index set \(I_c\) of Hilbert
polynomial labels is finite by hypothesis, the amplitude interval is
bounded by the two integers \(a_c\le b_c\), and \(N_c\in\mathbb Z_{\ge0}\)
is a finite integer.  The vanishing of the four defect ranks says that
these labels cover exactly the retained classes imported from the HN
type-bound and semistable-substack rows.  The argument constructs no
universal family over \(\mathfrak M^{ss}_{R,c}\), no closed inertia
stratification, no flag stack, no cosection atlas, no transition map,
no Pfaffian orientation, and no protected trace.
\end{proof}

The packet
\(\texttt{certificates/moduli/retained\_class\_bounds}\) checks this
row.  It imports the retained HN type-bound rows and retained
semistable substacks, records the finite class set \(C_R\), records the
Hilbert-polynomial labels \(P_{c,i}\), records the amplitude intervals
\([a_c,b_c]\), and records finite regularity bounds \(N_c\).  Its
verifier status is
\texttt{RETAINED\_CLASS\_BOUNDS\_VERIFIED}.  The mirrored class-bound
data are the Hilbert-polynomial, amplitude, and regularity columns in
\(\texttt{certificates/moduli/k3e\_finite\_moduli/hn\_type\_bounds.csv}\),
and they discharge items \(170\)--\(173\).  They do not construct
universal complexes, finite closed inertia stratifications, extension
or flag stacks, cosection atlases, transitions, compact Hall stages,
Pfaffian orientations, or protected traces.

\begin{definition}[Retained universal perfect-complex datum]
\label{def:retained-universal-perfect-complex-datum}
Fix the retained finite substacks
\(\mathfrak M^{ss}_{R,c}\), their scalar rigidifications, and the
retained class-bound datum.  A retained universal perfect-complex datum
assigns to every \(c\in C_R\) a perfect complex
\[
\mathcal U_{R,c}\in
\operatorname{Perf}\bigl(X\times\mathfrak M^{\mathrm{rig}}_{R,c}\bigr),
\]
a base-change identification for every test morphism
\[
T\longrightarrow \mathfrak M^{\mathrm{rig}}_{R,c},
\qquad
\mathcal U_{R,c}|_{X\times T}\simeq E_T,
\]
and a finite Tor-amplitude interval
\[
\operatorname{TorAmp}(\mathcal U_{R,c})\subset [a_c,b_c].
\]
It carries two defect ranks
\[
d^{\mathrm{desc}}_{R,c}=0,\qquad d^{\mathrm{perf}}_{R,c}=0.
\]
The first says that the tautological complex on the bounded
Quot/Postnikov chart descends across the scalar rigidification.  The
second says that the descended object remains perfect on the retained
finite substack.  This is a universal-complex datum.  It is not a
finite closed inertia stratification, not an extension or flag stack,
not a cosection atlas, not a transition system, and not an orientation
datum.
\end{definition}

\begin{proposition}[Universal perfect complexes on retained finite substacks]
\label{prop:retained-universal-perfect-complexes}
\textnormal{[proved]}
Let the retained finite substacks carry the datum of
Definition~\ref{def:retained-universal-perfect-complex-datum}.  Then
each retained scalar-rigidified finite substack
\(\mathfrak M^{\mathrm{rig}}_{R,c}\) carries a universal perfect
complex \(\mathcal U_{R,c}\) on \(X\times
\mathfrak M^{\mathrm{rig}}_{R,c}\), compatible with base change and of
finite Tor amplitude contained in \([a_c,b_c]\).
\end{proposition}

\begin{proof}
The bounded Quot/Postnikov ambient chart carries its tautological
complex.  Restricting it to the retained finite substack gives the
tautological complex before rigidification.  The equality
\(d^{\mathrm{desc}}_{R,c}=0\) is precisely the descent condition across
the scalar rigidification
\(\mathfrak M^{\mathrm{der}}_{R,c}\to
\mathfrak M^{\mathrm{rig}}_{R,c}\).  The equality
\(d^{\mathrm{perf}}_{R,c}=0\) says that the descended complex is
perfect on \(X\times\mathfrak M^{\mathrm{rig}}_{R,c}\).  The datum
also supplies the base-change identifications and the Tor-amplitude
containment in \([a_c,b_c]\).  Thus the retained substack carries the
asserted universal perfect complex.  The proof constructs no finite
closed inertia stratification, extension or flag stack, cosection
atlas, transition map, Pfaffian orientation, or protected trace.
\end{proof}

The packet
\(\texttt{certificates/moduli/retained\_universal\_complexes}\) checks
this row.  It imports the retained semistable substacks, class-bound
rows, derived enhancements, and scalar rigidifications, records one
universal perfect complex per retained finite substack, verifies
base-change compatibility, finite Tor amplitude, zero descent defect,
and zero perfection defect.  Its verifier status is
\texttt{RETAINED\_UNIVERSAL\_COMPLEXES\_VERIFIED}.  The mirrored rows
in
\(\texttt{certificates/moduli/k3e\_finite\_moduli/universal\_complexes.csv}\)
discharge item \(174\).  They do not construct finite closed inertia
stratifications, extension or flag stacks, cosection atlases,
transitions, compact Hall stages, Pfaffian orientations, or protected
traces.

\begin{definition}[Retained \(E\)-translation rigidification datum]
\label{def:retained-e-translation-rigidification-datum}
Fix the retained scalar-rigidified finite substacks
\(\mathfrak M^{\mathrm{rig}}_{R,c}\) and their universal perfect
complexes.  A retained \(E\)-translation rigidification datum assigns
to every \(c\in C_R\) an action
\[
a_{R,c}:E\times \mathfrak M^{\mathrm{rig}}_{R,c}
\longrightarrow \mathfrak M^{\mathrm{rig}}_{R,c}
\]
with identity element \(0_E\), free on the retained row, together with
the quotient stack
\[
\mathfrak M^{E\text{-}\mathrm{rig}}_{R,c}
=\bigl[\mathfrak M^{\mathrm{rig}}_{R,c}/E\bigr].
\]
The datum carries a defect rank
\[
d^{E\text{-}\mathrm{rig}}_{R,c}=0.
\]
Thus the elliptic-translation direction has been removed without
defect on the retained finite substack.  This is an
\(E\)-translation-rigidification datum.  It is not a finite closed
inertia stratification, not an extension or flag stack, not a
cosection atlas, not a transition system, and not an orientation datum.
\end{definition}

\begin{proposition}[\(E\)-translation rigidifications on retained finite substacks]
\label{prop:retained-e-translation-rigidifications}
\textnormal{[proved]}
Let the retained scalar-rigidified finite substacks carry the datum of
Definition~\ref{def:retained-e-translation-rigidification-datum}.  Then
each retained finite substack has an \(E\)-translation rigidification
\(\mathfrak M^{E\text{-}\mathrm{rig}}_{R,c}\), and the translation
direction has zero retained quotient defect.
\end{proposition}

\begin{proof}
The datum gives an action of the elliptic curve \(E\) on
\(\mathfrak M^{\mathrm{rig}}_{R,c}\), with identity \(0_E\), and states
that the action is free on the retained row.  The quotient
\[
\bigl[\mathfrak M^{\mathrm{rig}}_{R,c}/E\bigr]
\]
therefore removes exactly the \(E\)-translation direction.  The
vanishing \(d^{E\text{-}\mathrm{rig}}_{R,c}=0\) says that the quotient
has no retained translation-rigidification defect.  The proof
constructs no finite closed inertia stratification, extension or flag
stack, cosection atlas, transition map, Pfaffian orientation, or
protected trace.
\end{proof}

The packet
\(\texttt{certificates/moduli/retained\_e\_translation\_rigidifications}\)
checks this row.  It imports scalar rigidifications and universal
perfect complexes, records the \(E\)-translation action on each
retained finite substack, verifies freeness on the retained row,
constructs the quotient by \(E\), and checks zero
translation-rigidification defect.  Its verifier status is
\texttt{RETAINED\_E\_TRANSLATION\_RIGIDIFICATIONS\_VERIFIED}.  The
mirrored rows in
\(\texttt{certificates/moduli/k3e\_finite\_moduli/e\_translation\_rigidifications.csv}\)
discharge item \(176\).  They do not construct finite closed inertia
stratifications, extension or flag stacks, cosection atlases,
transitions, compact Hall stages, Pfaffian orientations, or protected
traces.

\begin{definition}[Retained closed-substack datum]
\label{def:retained-closed-substack-datum}
Fix the retained \(E\)-translation-rigidified finite substacks
\(\mathfrak M^{E\text{-}\mathrm{rig}}_{R,c}\).  A retained
closed-substack datum assigns to every \(c\in C_R\) a finite family of
closed substacks
\[
\mathfrak Z_{R,c,j}\hookrightarrow
\mathfrak M^{E\text{-}\mathrm{rig}}_{R,c},
\qquad 1\le j\le m_{R,c},
\]
which covers the retained row.  It carries the finite integers
\[
m_{R,c}>0
\]
and defect ranks
\[
d^{\mathrm{cov}}_{R,c}=0,\qquad d^{\mathrm{in}}_{R,c}=0.
\]
The first defect rank says that the closed substacks cover the
retained row.  The second says that the finite residual inertia from
the scalar rigidification remains finite on every retained closed
stratum.  This is a retained closed-substack datum.  It is not closure
under extensions, not closure under Harder--Narasimhan factors, not
closure under duals, not an extension or flag stack, not a cosection
atlas, not a transition system, and not an orientation datum.
\end{definition}

\begin{proposition}[Retained closed finite substacks]
\label{prop:retained-closed-substacks}
\textnormal{[proved]}
Let the retained \(E\)-translation-rigidified finite substacks carry
the datum of Definition~\ref{def:retained-closed-substack-datum}.  Then
each retained finite row has a finite closed cover by substacks on
which the residual inertia is finite.  Equivalently, the finite
stratification row has positive strata count, finite residual inertia,
zero coverage defect, and zero inertia defect.
\end{proposition}

\begin{proof}
The retained semistable rows are finite type and specialization-closed
by Proposition~\ref{prop:retained-finite-type-semistable-substacks}.
The scalar rigidification removes \(\mathbb G_m\) and leaves finite
residual inertia by
Proposition~\ref{prop:retained-finite-residual-inertia-after-rigidification}.
The \(E\)-translation rigidification removes the elliptic translation
direction by
Proposition~\ref{prop:retained-e-translation-rigidifications}.  The
closed-substack datum then supplies the closed immersions
\(\mathfrak Z_{R,c,j}\hookrightarrow
\mathfrak M^{E\text{-}\mathrm{rig}}_{R,c}\), their finite cover, and
the equalities \(d^{\mathrm{cov}}_{R,c}=d^{\mathrm{in}}_{R,c}=0\).
Thus the retained row is represented by a finite closed cover with
finite residual inertia.  The proof constructs no extension closure,
HN-factor closure, dual closure, flag stack, cosection atlas,
transition map, Pfaffian orientation, or protected trace.
\end{proof}

The packet
\(\texttt{certificates/moduli/retained\_closed\_substacks}\) checks
this row.  It imports the retained finite-type semistable substacks,
the scalar rigidifications with finite residual inertia, and the
\(E\)-translation rigidifications, records one retained closed substack
and one closed-cover row per retained finite row, and checks zero
coverage and inertia defects.  Its verifier status is
\texttt{RETAINED\_CLOSED\_SUBSTACKS\_VERIFIED}.  The mirrored rows in
\(\texttt{certificates/moduli/k3e\_finite\_moduli/stratifications.csv}\)
discharge item \(177\).  They do not prove closure under extensions,
closure under HN factors, closure under duals, extension or flag
stacks, cosection atlases, transitions, compact Hall stages, Pfaffian
orientations, or protected traces.

\begin{definition}[Retained extension-closure datum]
\label{def:retained-extension-closure-datum}
Fix the retained finite HN word set
\[
\mathcal W_R^{\mathrm{HN}}
=\{0,s_3,s_2,s_1,e_{32},e_{21},f_{321}\}
\]
from the bounded HN type datum.  A retained extension-closure datum is
a finite set of binary extension rows
\[
u\ast v\longmapsto w,\qquad u,v,w\in \mathcal W_R^{\mathrm{HN}},
\]
with \(w\) retained and with defect rank
\[
d^{\mathrm{ext}}_{u,v}=0.
\]
For the first retained window the nontrivial rows are
\[
s_3\ast s_2=e_{32},\qquad
s_2\ast s_1=e_{21},\qquad
e_{32}\ast s_1=f_{321},\qquad
s_3\ast e_{21}=f_{321}.
\]
This is an extension-closure datum in the retained exact window.  It
is not an extension stack, not a two-step flag stack, not closure under
Harder--Narasimhan factors, not closure under duals, not a cosection
atlas, not a transition system, and not an orientation datum.
\end{definition}

\begin{proposition}[Retained substacks are closed under extensions]
\label{prop:retained-extension-closure}
\textnormal{[proved]}
Let the retained finite HN window carry the datum of
Definition~\ref{def:retained-extension-closure-datum}.  Then every
extension represented by one of the retained binary extension rows has
middle term again in the retained finite HN window.  The extension
closure defect is zero on every such row.
\end{proposition}

\begin{proof}
The bounded HN type datum supplies the finite word set
\(\mathcal W_R^{\mathrm{HN}}\).  The extension-closure datum supplies
the four nonempty binary rows displayed in
Definition~\ref{def:retained-extension-closure-datum}; their middle
terms are \(e_{32}\), \(e_{21}\), and \(f_{321}\), all elements of
\(\mathcal W_R^{\mathrm{HN}}\).  The equalities
\[
d^{\mathrm{ext}}_{s_3,s_2}
=d^{\mathrm{ext}}_{s_2,s_1}
=d^{\mathrm{ext}}_{e_{32},s_1}
=d^{\mathrm{ext}}_{s_3,e_{21}}=0
\]
say that no retained extension leaves the finite window.  Thus the
retained HN rows are closed under the supplied extensions.  The proof
constructs no extension stack, two-step flag stack, HN-factor closure,
dual closure, cosection atlas, transition map, Pfaffian orientation, or
protected trace.
\end{proof}

The packet
\(\texttt{certificates/moduli/retained\_extension\_closure}\) checks
this row.  It imports the retained HN type-bound packet and retained
closed-substack packet, records four nonempty binary extension rows,
checks that every middle HN type lies inside the retained finite HN
word set, and checks zero extension-closure defect.  Its verifier
status is \texttt{RETAINED\_EXTENSION\_CLOSURE\_VERIFIED}.  The
mirrored rows in
\(\texttt{certificates/moduli/k3e\_finite\_moduli/extension\_closure.csv}\)
discharge item \(178\).  They do not construct extension or flag
stacks, prove closure under HN factors, prove closure under duals,
construct cosection atlases, construct transitions, construct compact
Hall stages, construct Pfaffian orientations, or construct protected
traces.

\begin{definition}[Retained HN-factor-closure datum]
\label{def:retained-hn-factor-closure-datum}
Fix the retained finite HN word set
\(\mathcal W_R^{\mathrm{HN}}\) and the retained semistable factor types
\[
\mathcal T_R=\{s_3,s_2,s_1\}.
\]
A retained HN-factor-closure datum assigns to every nonzero word
\[
w=(t_1,\ldots,t_k)\in \mathcal W_R^{\mathrm{HN}}
\]
its ordered factor list \(t_i\in\mathcal T_R\), together with the
retained class \(C_{R,t_i}\), the finite-type semistable substack
\(\mathfrak M^{ss}_{R,t_i}\), and the retained closed substack
\(\mathfrak Z_{R,t_i}\) for each factor.  It carries defect ranks
\[
d^{\mathrm{HNfac}}_{w,i}=0,\qquad 1\le i\le k.
\]
This is HN-factor closure in the retained finite window.  It is not
closure under duals, not an extension or flag stack, not subquotient
closure in compactified flag stacks, not a cosection atlas, not a
transition system, and not an orientation datum.
\end{definition}

\begin{proposition}[Retained substacks are closed under HN factors]
\label{prop:retained-hn-factor-closure}
\textnormal{[proved]}
Let the retained finite HN window carry the datum of
Definition~\ref{def:retained-hn-factor-closure-datum}.  Then every HN
factor of every retained finite HN word is represented by a retained
finite-type semistable substack and by a retained closed substack.  All
HN-factor-closure defects vanish.
\end{proposition}

\begin{proof}
The bounded HN type datum supplies the retained words
\[
0,\ s_3,\ s_2,\ s_1,\ e_{32}=(s_3,s_2),\
e_{21}=(s_2,s_1),\ f_{321}=(s_3,s_2,s_1).
\]
The datum assigns to each nonzero factor occurrence the corresponding
retained class, finite-type semistable substack, and retained closed
substack.  There are ten such occurrences:
three one-factor words, two two-factor words, and one three-factor
word.  The equalities \(d^{\mathrm{HNfac}}_{w,i}=0\) say that each
factor occurrence lies in the retained finite window.  Thus passing
from a retained object to its HN factors does not leave the retained
substack system.  The proof constructs no compactified extension
stack, no two-step flag stack, no subquotient closure theorem, no dual
closure, no cosection atlas, no transition map, no Pfaffian
orientation, and no protected trace.
\end{proof}

The packet
\(\texttt{certificates/moduli/retained\_hn\_factor\_closure}\) checks
this row.  It imports retained HN type bounds, class bounds,
finite-type semistable substacks, retained closed substacks, and
retained extension closure, records ten HN-factor occurrence rows, and
checks that every factor occurrence is represented by a retained class,
a retained finite-type semistable substack, and a retained closed
substack.  Its verifier status is
\texttt{RETAINED\_HN\_FACTOR\_CLOSURE\_VERIFIED}.  The mirrored rows in
\(\texttt{certificates/moduli/k3e\_finite\_moduli/hn\_factor\_closure.csv}\)
discharge item \(179\).  They are not the dual-closure rows below, do
not construct extension or flag stacks, do not prove subquotient
closure in flag stacks, do not construct cosection atlases, do not
construct transitions, do not construct compact Hall stages, do not
construct Pfaffian orientations, and do not construct protected traces.

\begin{definition}[Retained dual-closure datum]
\label{def:retained-dual-closure-datum}
Fix the retained finite HN word set
\(\mathcal W_R^{\mathrm{HN}}\).  A retained dual-closure datum is an
involution \(D_R\) on the retained semistable factor types
\[
D_R(s_3)=s_1,\qquad D_R(s_2)=s_2,\qquad D_R(s_1)=s_3,
\]
together with the induced involution on HN words
\[
D_R(t_1,\ldots,t_k)
=(D_R(t_k),\ldots,D_R(t_1)).
\]
For the first retained window this gives
\[
D_R(0)=0,\quad D_R(s_3)=s_1,\quad D_R(s_2)=s_2,\quad
D_R(s_1)=s_3,\quad D_R(e_{32})=e_{21},\quad
D_R(e_{21})=e_{32},\quad D_R(f_{321})=f_{321}.
\]
It carries defect ranks
\[
d^{\vee}_{w}=0,\qquad w\in\mathcal W_R^{\mathrm{HN}}.
\]
This is dual closure in the retained finite window.  It is not an
extension or flag stack, not subquotient closure in compactified flag
stacks, not a cosection atlas, not a transition system, and not an
orientation datum.
\end{definition}

\begin{proposition}[Retained substacks are closed under duals]
\label{prop:retained-dual-closure}
\textnormal{[proved]}
Let the retained finite HN window carry the datum of
Definition~\ref{def:retained-dual-closure-datum}.  Then duality sends
every retained finite HN word to a retained finite HN word, and the
dual-closure defect vanishes on every retained word.
\end{proposition}

\begin{proof}
The factor duality exchanges \(s_3\) and \(s_1\) and fixes \(s_2\), so
it is an involution on the retained semistable factor set
\(\{s_3,s_2,s_1\}\).  Applying it to an HN word reverses the order and
dualizes each factor.  Hence
\[
D_R(e_{32})=D_R(s_3,s_2)=(s_2,s_1)=e_{21},\qquad
D_R(e_{21})=(s_3,s_2)=e_{32},
\]
and
\[
D_R(f_{321})=D_R(s_3,s_2,s_1)=(s_3,s_2,s_1)=f_{321}.
\]
The zero word is fixed, the one-factor words map as displayed in the
definition, and all targets belong to \(\mathcal W_R^{\mathrm{HN}}\).
The equalities \(d^{\vee}_w=0\) say that no dualized retained word
leaves the finite window.  The proof constructs no compactified
extension stack, no two-step flag stack, no subquotient closure
theorem, no cosection atlas, no transition map, no Pfaffian
orientation, and no protected trace.
\end{proof}

The packet
\(\texttt{certificates/moduli/retained\_dual\_closure}\) checks this
row.  It imports retained HN type bounds, class bounds, retained closed
substacks, and retained HN-factor closure, records the three factor
dual rows and seven HN-word dual rows, verifies that type duality and
HN-word duality are involutive, and checks that HN-word duality is
reverse factorwise duality.  Its verifier status is
\texttt{RETAINED\_DUAL\_CLOSURE\_VERIFIED}.  The mirrored rows in
\(\texttt{certificates/moduli/k3e\_finite\_moduli/dual\_closure.csv}\)
discharge item \(180\).  They do not construct extension or flag
stacks, prove subquotient closure in flag stacks, construct cosection
atlases, construct transitions, construct compact Hall stages,
construct Pfaffian orientations, or construct protected traces.

\begin{definition}[Algebraic/effective Hall support]
\label{def:algebraic-effective-hall-support}
Fix a polarised K3 surface $(S, H)$, a chosen Mukai vector lattice
$\Lambda_{S, \mathrm{alg}} \subset \Lambda_S$ (the image of the
algebraic K-class lattice generated by sheaves whose Mukai vector lies in
$H^0\oplus\mathrm{NS}(S)\oplus H^4$), and the retained stability datum
\(\sigma=\sigma_{s,t}\) with heart
\(\mathcal A_X^{\mathrm{AP/Liu}}\) from
Definition~\ref{def:retained-ap-liu-heart}.  The retained-window rows
above supply noetherianity, Harder--Narasimhan filtrations, bounded HN
types, finite-type semistable substacks, and quasi-smooth derived
enhancements, together with finite residual inertia after scalar
rigidification, finite retained class-bound data, and universal
perfect complexes, \(E\)-translation rigidifications, and finite
closed-cover stratifications at finite height, together with extension
closure, HN-factor closure, and dual closure inside the retained HN
window; later rows supply flag stacks, subquotient closure in flag
stacks, cosection charts, transitions, scalar firewalls, and
compact Hall support.  No unconditional Bridgeland stability theorem on the
threefold \(S\times E\) is used.  Set
\[
\Gamma_X^{\mathrm{alg}}
:= \Lambda_{S,\mathrm{alg}}\oplus\Lambda_{S,\mathrm{alg}},\qquad
\Gamma_{X, \sigma}^{\mathrm{eff}}
\subset \Gamma_X^{\mathrm{alg}}
\subset \Gamma_X^{\mathrm{form}},
\]
where the \emph{$\sigma$-effective semigroup}
$\Gamma_{X, \sigma}^{\mathrm{eff}}$ is the additive sub-semigroup of
charges supported by objects semistable for this datum and of nonzero
Mukai vector.
Membership in $\Gamma_{X, \sigma}^{\mathrm{eff}}$ is not implied by
membership in $\Gamma_X^{\mathrm{alg}}$ or by formal lattice arithmetic
(folklore).
\end{definition}

\begin{definition}[Admissible Dirac--Igusa parameter datum]
\label{def:admissible-dirac-igusa-parameter-datum}
An admissible parameter datum for the Dirac--Igusa construction is
\[
\mathfrak p=(S,H,\sigma,E,0_E,\tau_X,\mathcal R_{\mathrm{HN}}),
\qquad X=S\times E,
\]
where \((S,H)\) is a polarized smooth complex K3 surface, \(E\) is a
smooth elliptic curve with origin \(0_E\), \(\tau_X:K_X\simeq\mathcal O_X\)
is the chosen trivialization, \(\sigma\) is the retained stability datum
used to define \(\Gamma_{X,\sigma}^{\mathrm{eff}}\), and
\(\mathcal R_{\mathrm{HN}}\) is the cofinal finite HN index system.
It includes the retained finite charge, moduli, hybrid, orientation,
Hall, Pfaffian, Koszul, and primitive-recognition packets whenever
those packets are invoked later.  Thus \(H\) and \(\sigma\) are not
decorations: they determine the algebraic/effective Hall support and
the finite HN windows.

A morphism \(\alpha:\mathfrak p\to\mathfrak p'\) is an admissible
base-change isomorphism
\[
\alpha=(f,u,\eta_\sigma,\eta_{\mathrm{HN}})
\]
with \(f:S\to S'\) an isomorphism of K3 surfaces carrying \(H\) to
\(H'\), \(u:E\to E'\) an isomorphism of elliptic curves with
\(u(0_E)=0_{E'}\), compatibility with the chosen trivializations of
\(K_{S\times E}\), a transport \(\eta_\sigma\) of the retained
stability datum and effective Hall support, and a cofinal identification
\(\eta_{\mathrm{HN}}:\mathcal R_{\mathrm{HN}}\to\mathcal R'_{\mathrm{HN}}\).
The induced Mukai isometry sends \((Q,P)\) to
\((f_*Q,f_*P)\), preserves the Mukai pairing, and therefore preserves
\[
\Pi_X(Q,P)=((Q,Q)/2,(Q,P),(P,P)/2),\qquad B=-\delta\Pi_X .
\]
The resulting groupoid is denoted \(\mathsf{Par}^{\mathrm{DI}}\).
Arbitrary non-isomorphic deformations of \(S\), arbitrary changes of
stability chamber, and arbitrary changes of elliptic origin are not
morphisms of \(\mathsf{Par}^{\mathrm{DI}}\) unless the corresponding
wall-crossing or origin-translation datum below is supplied.
\end{definition}

\begin{proposition}[Admissible base-change functoriality]
\label{prop:admissible-base-change-functoriality}
Let \(\alpha:\mathfrak p\to\mathfrak p'\) be a morphism in
\(\mathsf{Par}^{\mathrm{DI}}\).  If the finite packets attached to
\(\mathfrak p\) have zero defect rows and are transported by
\(\alpha\), then pull-push along \(f\times u:X\to X'\) induces
functors at each retained height
\[
\alpha_{R,*}:\mathsf{DI}^{\mathrm{fin}}_{X,R}
\longrightarrow
\mathsf{DI}^{\mathrm{fin}}_{X',\eta_{\mathrm{HN}}(R)}
\]
preserving the components \(\mathcal A^{E_3}\), \(F^{\mathrm{hyb}}\),
\(\widehat\Gamma\), \(\Pi\), \(\Phi\), \(o\), \(H\), \(P^\Pi\), \(C\),
\(\Theta_{\mathrm{Kos}}\), \(\mathscr L_{\mathrm{Pf}}\),
\(\operatorname{pf}\), \(\varepsilon_o\), and \(\mathrm{Rec}\), in the
sense of the morphisms \textup{(M1)}--\textup{(M8)} of
Definition~\ref{def:finite-stage-dirac-igusa-morphism}.  These functors
commute with HN transition morphisms and therefore induce a functor on
Dirac--Igusa pro-objects
\[
\alpha_*:
\operatorname{Pro}(\mathsf{DI}^{\mathrm{fin}}_X)
\longrightarrow
\operatorname{Pro}(\mathsf{DI}^{\mathrm{fin}}_{X'}).
\]
For composable admissible morphisms, the induced functors compose up to
the canonical pro-isomorphism coming from the common cofinal HN
refinement.  This is the base-change functoriality in \(S\) and \(E\)
used by the construction.
\end{proposition}

\begin{proof}
The isomorphism \(f\) gives an isometry of Mukai lattices and carries
the algebraic sublattice, the polarization, and the retained stability
support by hypothesis.  The isomorphism \(u\) preserves \(1_E\),
\(\omega_E\), the origin, and the translation group law.  Hence
\((f\times u)_*\) preserves the formal charge lattice, the
normal-ordered central extension, and the Gram map.  The assumed
transport of the finite packets carries the moduli stacks,
correspondence stacks, orientation lines, Pfaffian lines, Hall
bialgebra matrices, Koszul source complexes, and primitive-recognition
matrices to their primed counterparts.  Compatibility with the
transition maps is exactly the zero-defect condition in the finite
morphism tables.  Proposition~\ref{prop:dirac-igusa-pro-object-universal-property}
then gives the induced pro-functor and the cofinal composition law.
\end{proof}

\begin{definition}[Polarization and elliptic-origin dependence]
\label{def:polarization-origin-dependence}
The dependence on the K3 polarization is the dependence of
\(\Gamma_{X,\sigma}^{\mathrm{eff}}\), the HN index system, and the
finite moduli stacks on \((H,\sigma)\).  If two choices
\((H,\sigma)\) and \((H',\sigma')\) lie in chambers connected by a
specified Kontsevich--Soibelman wall-crossing datum
\(\mathsf{WC}_{H,H',R}\) with zero finite defects, the resulting
Dirac--Igusa pro-objects are compared by the induced pro-isomorphism.
Without such a wall-crossing datum they are different parameter
points, not two presentations of the same object.

Let \(a\in E\) and replace the origin \(0_E\) by \(a\).  For a retained
wrapped family \(\mathcal F\), the determinant anchors of
Definition~\ref{def:determinant-anchor} satisfy the canonical identity
\[
\lambda^{\det}_{a}(\mathcal F)
\simeq
\lambda^{\det}_{0_E}(\mathcal F)\otimes
\mathcal O_E\!\bigl(\chi(\mathcal F)(0_E-a)\bigr)
\quad\text{in }\operatorname{Pic}^0(E).
\]
Thus changing the origin contributes the degree-zero character
\[
\vartheta_{a,\mathcal F}:=
\mathcal O_E\!\bigl(\chi(\mathcal F)(0_E-a)\bigr)\in \operatorname{Pic}^0(E),
\]
and, after restriction to a finite stabilizer \(H\subset E[N]\), the
corresponding finite character
\[
\vartheta_{a,\mathcal F}^{H}\in H^1(BH;\mathbb C^\times)
\]
is the only permitted change in the quotient-orientation linearization.
The Dirac--Igusa object is independent of the elliptic origin only
after twisting the wrapped-anchor and orientation-linearization rows by
this character and verifying the same finite stabilizer equations.
\end{definition}

\begin{proposition}[Elliptic-origin change up to its canonical character]
\label{prop:elliptic-origin-independence-character}
Fix \(S,H,\sigma,E\) and the finite HN system, and let
\(\mathfrak p_0\) and \(\mathfrak p_a\) be the admissible parameter
data obtained from the origins \(0_E\) and \(a\in E\).  Suppose the
finite hybrid, orientation, and morphism rows for \(\mathfrak p_a\) are
obtained from those for \(\mathfrak p_0\) by tensoring every wrapped
anchor with
\(\vartheta_{a,\mathcal F}=\mathcal O_E(\chi(\mathcal F)(0_E-a))\)
and by adding the restricted character
\(\vartheta^H_{a,\mathcal F}\) to the finite stabilizer linearization
row for each retained stabilizer \(H\subset E[N]\).  Then the two
finite-stage systems determine canonically pro-isomorphic
Dirac--Igusa pro-objects after this character twist.  Before the twist,
the obstruction to identifying the two origin choices is exactly the
family of characters \(\vartheta^H_{a,\mathcal F}\).
\end{proposition}

\begin{proof}
By Definition~\ref{def:determinant-anchor},
\[
\lambda^{\det}_{a}(\mathcal F)
=\det Rp_{E,*}\mathcal F\otimes\mathcal O_E(-\chi(\mathcal F)a)
\]
and
\[
\lambda^{\det}_{0_E}(\mathcal F)
=\det Rp_{E,*}\mathcal F\otimes\mathcal O_E(-\chi(\mathcal F)0_E).
\]
Their ratio is
\(\mathcal O_E(\chi(\mathcal F)(0_E-a))\), a canonical point of
\(\operatorname{Pic}^0(E)\).  Restriction to a finite stabilizer gives
the character \(\vartheta^H_{a,\mathcal F}\).  All other components of
the formal charge lattice, the Mukai--Gram map, and the K3 data are
unchanged.  Therefore the finite-stage morphisms are the identity on
all components except the wrapped-anchor and orientation-linearization
rows, where they are precisely the displayed tensor and character
twists.  The zero-defect assumptions give commuting transition squares;
Proposition~\ref{prop:dirac-igusa-pro-object-universal-property} then
identifies the resulting pro-objects.  If the character rows are not
twisted, the displayed ratio remains as the residual finite-stabilizer
linearization class, so the obstruction is exactly
\(\vartheta^H_{a,\mathcal F}\).
\end{proof}

\begin{remark}[Three-tier inclusion]
\label{rem:three-tier-inclusion}
The hierarchy $\Gamma_{X, \sigma}^{\mathrm{eff}}
\subset \Gamma_X^{\mathrm{alg}}
\subset \Gamma_X^{\mathrm{form}}$
distinguishes three discipline registers.  $\Gamma_X^{\mathrm{form}}$ is
the formal additive group on which $\Pi_X$ is defined and the
polarization identity $B = -\delta\Pi_X$ holds.  $\Gamma_X^{\mathrm{alg}}$
is the K-theoretic algebraic sublattice; its rank-one extension closure
inside $\Gamma_X^{\mathrm{form}}$ records only the algebraic Mukai
branch (computed, \cite{Mukai1987,KSCoHA}).
$\Gamma_{X, \sigma}^{\mathrm{eff}}$ is the Hall support selected by
$\sigma$; effectivity and semistability are not lattice-theoretic
properties and are not fibre-functorial in the K3 type.  The sectorial
Hall truncation of \S\ref{sec:k3xe-hybrid-ran-carrier} works at this
third level after imposing finite Harder--Narasimhan height bounds.
\end{remark}

\begin{lemma}[Two orthogonal hyperbolic planes inside $\Lambda_S$]
\label{lem:two-hyperbolic-planes}
The Mukai lattice contains two orthogonal hyperbolic planes:
\[
U \oplus U \hookrightarrow \Lambda_S,\qquad
U_1 = \Z e_1 \oplus \Z f_1,\quad U_2 = \Z e_2 \oplus \Z f_2,
\]
with $e_i^2 = f_i^2 = 0$, $e_i\cdot f_i = 1$, and $U_1\perp U_2$ inside
the Mukai pairing (proved, \cite{NIKULIN}).
\end{lemma}

\begin{proof}
Direct from $\Lambda_S \cong U^{\oplus 4} \oplus E_8(-1)^{\oplus 2}$.
\end{proof}

\section{Quadratic Gram map and lattice polarization}
\label{sec:gram-map-polarization}

\subsection{Igusa variables and the Gram-index lattice}
\label{subsec:igusa-variables-gram-index}

For a Siegel period
\[
Z = \begin{pmatrix} z_1 & z_2 \\ z_2 & z_3 \end{pmatrix}
\in \mathcal H_2,\qquad
q = e^{2\pi i z_1},\quad r = e^{2\pi i z_2},\quad s = e^{2\pi i z_3},
\]
and
\[
T(n, l, m) = \begin{pmatrix} n & l/2 \\ l/2 & m \end{pmatrix},\qquad
n, l, m \in \Z,
\]
the trace pairing
\[
\bigl(T(n, l, m), Z\bigr)_{\mathrm{tr}}
:= \mathrm{tr}\bigl(T(n, l, m)\,Z\bigr)
= n z_1 + l z_2 + m z_3
\]
gives the Borcherds monomial $x^{(n, l, m)} = q^n r^l s^m$.

\begin{definition}[Gram-index lattice and effective semigroup]
\label{def:gram-index-lattice}
The \emph{Gram-index lattice} is
\[
\Gamma_{\mathrm{gram}} := \Z^3,\qquad
\gamma = (n, l, m).
\]
The \emph{cusp-positive semigroup} at the standard cusp $\mathfrak
c_\infty$ with chamber convention $0 < |s| \ll |q| \ll |r|^{-1} \ll 1$
is
\[
\Gamma_{\mathrm{eff}}
= \{(n, l, m) : m > 0,\ n \ge 0,\ l \in \Z\}
\cup \{(n, l, 0) : n > 0,\ l \in \Z\}
\cup \{(0, l, 0) : l < 0\}.
\]
The \emph{Jacobi-active sublocus}
$\Gamma_{\mathrm{act}} \subset \Gamma_{\mathrm{eff}}$ is the locus of
nonvanishing $f(nm, l)$ in the weak Jacobi form $\phi_{0, 1}$
(computed, \cite{EZ:1}).
\end{definition}

\begin{proposition}[Gram matrix convention for protected charges]
\label{prop:gram-matrix-convention}
For a formal even charge $c = (Q, P) \in \Gamma_X^{\mathrm{form}}$ in
the unimodular lattice $\Lambda_S$, set
\[
n = \frac{(Q, Q)}{2},\qquad
l = (Q, P),\qquad
m = \frac{(P, P)}{2}.
\]
Then $n, l, m \in \Z$ (since $\Lambda_S$ is even unimodular,
proved),
\[
\begin{pmatrix} (Q, Q) & (Q, P) \\ (P, Q) & (P, P) \end{pmatrix}
= 2\,T(n, l, m),\qquad
\det\bigl(2 T(n, l, m)\bigr) = 4 n m - l^2,
\]
and $x^{(n, l, m)} = q^n r^l s^m$ is the Fourier character of
$\bigl(T(n, l, m), Z\bigr)_{\mathrm{tr}}$.
\end{proposition}

\begin{proof}
The Mukai pairing is integer-valued and even on $\Lambda_S\oplus\Lambda_S$
by Definition~\ref{def:mukai-lattice}.  The matrix identity is the
direct expansion
\[
\begin{pmatrix} (Q, Q) & (Q, P) \\ (P, Q) & (P, P) \end{pmatrix}
= \begin{pmatrix} 2 n & l \\ l & 2 m \end{pmatrix}
= 2\,T(n, l, m).
\]
The determinant follows.  The Fourier character is
$\exp(2\pi i\,\mathrm{tr}(T Z)) = q^n r^l s^m$.
\end{proof}

\subsection{The Gram map and its bilinear cocycle}
\label{subsec:gram-map-bilinear-cocycle}

\begin{definition}[Quadratic Gram map and bilinear cocycle]
\label{def:gram-map-cocycle}
The \emph{Gram map} is the quadratic projection
\[
\Pi_X : \Gamma_X^{\mathrm{form}} \longrightarrow \Gamma_{\mathrm{gram}},
\qquad
\Pi_X(Q, P)
= \left(\frac{(Q, Q)}{2},\,(Q, P),\,\frac{(P, P)}{2}\right).
\]
The \emph{symmetric bilinear cocycle} $B$ associated to $\Pi_X$ is
\[
B : \Gamma_X^{\mathrm{form}} \times \Gamma_X^{\mathrm{form}}
\longrightarrow \Gamma_{\mathrm{gram}},\qquad
B\bigl((Q, P), (Q', P')\bigr)
= \bigl((Q, Q'),\,(Q, P') + (Q', P),\,(P, P')\bigr).
\]
\end{definition}

\begin{lemma}[Gram additivity defect]
\label{lem:gram-additivity-defect}
For $c, c' \in \Gamma_X^{\mathrm{form}}$,
\[
\Pi_X(c + c') = \Pi_X(c) + \Pi_X(c') + B(c, c').
\]
The defect $B$ is symmetric, $\Z$-bilinear, integer-valued, and
satisfies polarization: $B(c, c) = 2\,\Pi_X(c)$.  As an ordinary
additive cochain identity,
\[
B = -\delta\Pi_X,\qquad
(\delta q)(c, c') = q(c) + q(c') - q(c + c').
\]
(proved.)
\end{lemma}

\begin{remark}[Quadratic primitive, no linear primitive]
\label{rem:B-quadratic-not-linear}
The equality \(B=-\delta\Pi_X\) is an equality in the ordinary
inhomogeneous cochain complex, where the primitive \(\Pi_X\) is
quadratic.  There is no additive homomorphism
\(L:\Gamma_X^{\mathrm{form}}\to\Gamma_{\mathrm{gram}}\) with
\(B=-\delta L\).  Thus the normal-ordering defect is removed by the
quadratic Gram coordinate, not by a linear change of charge lattice.
\end{remark}

\begin{proof}
Write $c = (Q, P)$, $c' = (Q', P')$; then
\[
\frac{(Q + Q')^2}{2}
= \frac{Q^2}{2} + \frac{Q'^2}{2} + Q\cdot Q',
\]
\[
(Q + Q')\cdot(P + P')
= Q\cdot P + Q'\cdot P' + Q\cdot P' + Q'\cdot P,
\]
\[
\frac{(P + P')^2}{2}
= \frac{P^2}{2} + \frac{P'^2}{2} + P\cdot P'.
\]
The three cross terms are the three components of $B(c, c')$.
Symmetry, bilinearity, and integrality follow from the corresponding
properties of the Mukai pairing.  Polarization is the case $c' = c$.
The cochain identity rewrites the additivity defect.
\end{proof}

\begin{lemma}[Bilinear cocycle identity]
\label{lem:b-cocycle-identity}
$B$ is normalized, $B(c, 0) = B(0, c) = 0$, and the additive coboundary
\[
(\delta B)(a, b, c) := B(b, c) - B(a + b, c) + B(a, b + c) - B(a, b)
\]
vanishes identically on $\Gamma_X^{\mathrm{form}}$ (proved).
\end{lemma}

\begin{proof}
By bilinearity in each slot,
\[
B(a + b, c) = B(a, c) + B(b, c),\qquad
B(a, b + c) = B(a, b) + B(a, c).
\]
Substituting,
\[
(\delta B)(a, b, c)
= B(b, c) - B(a, c) - B(b, c) + B(a, b) + B(a, c) - B(a, b) = 0.
\]
The polarization identity $B(c, c) = 2\,\Pi_X(c)$ is the only place where
$B$ takes a non-vanishing value on the diagonal; this is structural,
not a Bott-element framing.
\end{proof}

\begin{remark}[Polarization, not Bott element]
\label{rem:polarization-not-bott}
The identity $B(c, c) = 2\,\Pi_X(c)$ is the polarization identity for
the symmetric bilinear $B$ associated to the quadratic $\Pi_X$.  It is
forced by the homogeneous-of-degree-two structure of the Gram map.  No
external Bott-element framing or characteristic-class coincidence is
involved.  The full structural content is that $\Pi_X$ is a homogeneous
quadratic on the additive group $\Gamma_X^{\mathrm{form}}$ valued in
$\Gamma_{\mathrm{gram}}$, with associated symmetric bilinear $B$ given by
the Mukai pairing of components on each direction.
\end{remark}

\subsection{Formal lift of every Gram triple}
\label{subsec:formal-gram-triple-lift}

\begin{lemma}[Formal primitive torsion-one lift of every Gram triple]
\label{lem:primitive-lift-every-gram-triple}
Suppose $\Lambda_S$ contains two orthogonal hyperbolic planes $U_1\oplus
U_2$ with bases $(e_1, f_1, e_2, f_2)$ as in
Lemma~\ref{lem:two-hyperbolic-planes}.  Then every Gram triple
$(n, l, m) \in \Z^3$ admits a primitive saturated formal full-Mukai lift
$c = (Q, P) \in \Gamma_X^{\mathrm{form}}$ with $\Pi_X(Q, P) = (n, l, m)$.
One may take
\[
Q = e_1 + n f_1,\qquad
P = l f_1 + e_2 + m f_2.
\]
For this lift, the wedge torsion invariant
\[
I(Q, P)
:= \gcd\!\bigl\{\text{coordinates of }Q\wedge P\text{ in a }\Z\text{-basis of }
{\textstyle\bigwedge^2\Lambda_S}\bigr\}
\]
equals $1$.
(proved.)
\end{lemma}

\begin{proof}
The hyperbolic-plane pairings give
\[
Q^2 = 2 n,\qquad P^2 = 2 m,\qquad Q\cdot P = l,
\]
hence $\Pi_X(Q, P) = (n, l, m)$.  The submodule $\Z Q + \Z P \subset
\Lambda_S$ contains the unimodular minor
$\det\!\begin{pmatrix} 1 & 0 \\ 0 & 1 \end{pmatrix} = 1$ in the rows
indexed by $e_1, e_2$ of the coefficient matrix in
$(e_1, f_1, e_2, f_2)$, so the lift is primitive saturated.
The wedge expansion
\[
Q\wedge P
= l\,e_1\wedge f_1 + e_1\wedge e_2 + m\,e_1\wedge f_2
+ n\,f_1\wedge e_2 + nm\,f_1\wedge f_2
\]
contains the basis element $e_1\wedge e_2$ with coefficient $1$, so the
gcd is $1$ and $I(Q, P) = 1$.
\end{proof}

\begin{remark}[Pointwise existence is not Hall support]
\label{rem:pointwise-not-hall-support}
Lemma~\ref{lem:primitive-lift-every-gram-triple} produces a formal
charge plane realizing every prescribed Igusa degree.  It does not prove
algebraicity of $(Q, P)$, effectivity, non-empty compact Hall moduli,
duality-orbit descent, source-index descent, or any chain-level Hall
strictification.  The compact Hall obstruction is functorial
compatibility of pointwise existence with Hall addition and the
primitive bracket; only the latter is the working content of
Chapter~\ref{chap:dirac-igusa-pro-object}.
\end{remark}

The packet
\(\texttt{certificates/charge/formal\_primitive\_lift}\) is the finite
table-level check of Lemma~\ref{lem:primitive-lift-every-gram-triple}.
It verifies, on the zero row, the three simple-root rows, the
\(\delta_{123}\) row, and one negative-pairing test row, that
\[
\Pi_X(Q,P)=(n,l,m),\qquad
I(Q,P)=1,\qquad
\alpha(\Pi_X(Q,P))=2n f_2-l f_3+2m f_{-2}.
\]
The same packet records the \(e_1,e_2\) minor equal to \(1\), hence
primitivity of the displayed formal lift.  It is not evidence for
algebraic effectivity, finite HN boundedness, compact Hall support,
source Hall brackets, Pfaffian orientations, type-II wall atlases,
mirror discriminants, or protected traces.

\section{Normal-ordered Gram extension}
\label{sec:normal-ordered-extension}

\subsection{Central extension and the additive normal-ordered Gram
homomorphism}
\label{subsec:central-extension-additive-gram}

\begin{definition}[Normal-ordered extension and Gram homomorphism]
\label{def:normal-ordered-extension}
Set
\[
\widehat\Gamma_X
:= \Gamma_X^{\mathrm{form}}\oplus\Gamma_{\mathrm{gram}},
\]
with operation
\[
(c, T)\star(c', T')
:= \bigl(c + c',\,T + T' + B(c, c')\bigr).
\]
Define the \emph{normal-ordered Gram map}
\[
\overline\Pi_X : \widehat\Gamma_X \longrightarrow \Gamma_{\mathrm{gram}},
\qquad
\overline\Pi_X(c, T) = \Pi_X(c) - T.
\]
The \emph{raw placement} and the \emph{split section} are
\[
i_0 : \Gamma_X^{\mathrm{form}}\longrightarrow\widehat\Gamma_X,\quad
i_0(c) = (c, 0),\qquad
s : \Gamma_X^{\mathrm{form}}\longrightarrow\widehat\Gamma_X,\quad
s(c) = (c, \Pi_X(c)).
\]
\end{definition}

\begin{theorem}[Group law of the normal-ordered extension; additivity of
$\overline\Pi_X$]
\label{thm:normal-ordered-extension-group-law}
Let $\widehat\Gamma_X$, $\star$, $\overline\Pi_X$, $i_0$, $s$ be as in
Definition~\ref{def:normal-ordered-extension}.  Then:
\begin{enumerate}[label = \textup{(\roman*)}]
\item $(\widehat\Gamma_X, \star)$ is an abelian group with identity
$(0, 0)$ and inverse $(c, T)^{-1} = (-c,\,-T + B(c, c)) = (-c,\,-T +
2\,\Pi_X(c))$.
\item The sequence
\[
0\to\Gamma_{\mathrm{gram}}\xrightarrow{T\mapsto(0, T)}\widehat\Gamma_X
\xrightarrow{\mathrm{pr}_1}\Gamma_X^{\mathrm{form}}\to 0
\]
is a central extension of abelian groups with cocycle $B$.
\item $\overline\Pi_X$ is a group homomorphism from
$(\widehat\Gamma_X, \star)$ to the additive group $\Gamma_{\mathrm{gram}}$.
\item The split section $s(c) = (c, \Pi_X(c))$ is additive, and
\[
\overline\Pi_X(i_0(c)) = \Pi_X(c),\qquad
\overline\Pi_X(s(c)) = 0.
\]
\item The change of variables $(c, T)\mapsto(c, T - \Pi_X(c))$ carries
$\star$ to componentwise addition; thus $\widehat\Gamma_X\cong
\Gamma_X^{\mathrm{form}}\oplus\Gamma_{\mathrm{gram}}$ as abstract abelian
groups, and $[B] = 0$ in additive cohomology.
\end{enumerate}
(proved.)
\end{theorem}

\begin{proof}
\emph{(i)} Associativity follows from the cocycle identity
$\delta B = 0$ (Lemma~\ref{lem:b-cocycle-identity}); the associator
$B(a, b) + B(a + b, c) - B(a, b + c) - B(b, c)$ vanishes by bilinearity.
Commutativity is symmetry of $B$.  The unit calculation $(c, T)\star(0,
0) = (c, T + B(c, 0)) = (c, T)$ and inverse $(c, T)\star(-c, -T + B(c,
c)) = (0, B(c, c) - B(c, c)) = (0, 0)$ (using $B(c, -c) = -B(c, c)$ by
bilinearity) close the abelian-group structure.

\emph{(ii)} Definition.

\emph{(iii)} Apply Lemma~\ref{lem:gram-additivity-defect}:
\[
\overline\Pi_X((c, T)\star(c', T'))
= \Pi_X(c + c') - T - T' - B(c, c')
= \Pi_X(c) + \Pi_X(c') - T - T'
= \overline\Pi_X(c, T) + \overline\Pi_X(c', T').
\]

\emph{(iv)} For the section,
\[
s(c)\star s(c')
= (c + c',\,\Pi_X(c) + \Pi_X(c') + B(c, c'))
= (c + c',\,\Pi_X(c + c'))
= s(c + c').
\]
The identities $\overline\Pi_X(i_0(c)) = \Pi_X(c) - 0 = \Pi_X(c)$ and
$\overline\Pi_X(s(c)) = \Pi_X(c) - \Pi_X(c) = 0$ are immediate.

\emph{(v)} The map $(c, T)\mapsto(c, T - \Pi_X(c))$ is a bijection with
inverse $(c, T)\mapsto(c, T + \Pi_X(c))$.  Direct calculation gives
\[
(c, T - \Pi_X(c)) + (c', T' - \Pi_X(c'))
= (c + c',\,T + T' - \Pi_X(c) - \Pi_X(c'))
= (c + c',\,(T + T' + B(c, c')) - \Pi_X(c + c')),
\]
matching the change of variables applied to $(c, T)\star(c', T')$.
\end{proof}

\begin{remark}[Finite charge-window verification]
\label{rem:finite-charge-window-verification}
The normal-ordered group law is a proved lattice identity.  Its use in
the Dirac--Igusa pro-object additionally requires a finite
charge-window packet: active \(\phi_{0,1}\)-support, downward
saturation, Liu HN bounds, Mukai--Gram rows, cocycle identities,
normal-ordered lift orbits, target-window images, primitive formal
lifts, strict transitions, and a scalar firewall.  A target
presentation, signed exponent list, Pfaffian product, Hilbert-scheme
scalar, or source-rank table does not by itself supply this packet.
\end{remark}

\subsection{The fixed-lift raw-pushforward obstruction theorem}
\label{subsec:raw-gram-descent-no-go}

\begin{definition}[Raw homogeneous and fixed-lift $\Pi_X$-pushforwards]
\label{def:raw-pi-pushforward}
Let $P = \bigoplus_{c\in\Gamma_X^{\mathrm{form}}} P_c$ be a
$\Gamma_X^{\mathrm{form}}$-graded primitive object.  The
\emph{raw homogeneous $\Pi_X$-pushforward} is
\[
\bigl(\Pi_{X, *}^{\mathrm{raw}} P\bigr)_\gamma
:= \bigoplus_{\Pi_X(c) = \gamma} P_c.
\]
The \emph{strict fixed-lift raw $\Pi_X$-pushforward} is the narrower
operation obtained by choosing one lift $L(\gamma) \in
\Gamma_X^{\mathrm{form}}$ for every retained $\gamma$ and setting
\[
\bigl(\Pi_{X, *}^{\mathrm{raw}, L} P\bigr)_\gamma
:= P_{L(\gamma)}.
\]
\end{definition}

\begin{theorem}[Raw homogeneous fixed-lift $\Pi_X$-pushforward cannot
realize the type-II real-root strings]
\label{thm:raw-gram-descent-no-go}
Let
$P = \bigoplus_{c\in\Gamma_X^{\mathrm{form}}} P_c$ be a
$\Gamma_X^{\mathrm{form}}$-graded Lie superalgebra with $[P_c, P_{c'}]
\subset P_{c + c'}$.  Suppose a strict fixed-lift raw pushforward of
$P$ along the quadratic map $\Pi_X$ were a $\Gamma_{\mathrm{gram}}$-graded
Lie superalgebra structure on the displayed bracket channels of chosen
charge lifts.  Then every nonzero bracket channel $[P_c, P_{c'}]$
contributing to that strict raw pushforward must satisfy $B(c, c') = 0$.

For the type-II real simple roots $(\delta_1, \delta_2, \delta_3)$ of
the Gritsenko--Nikulin algebra $\mathfrak g_{\Delta_5}$ with Cartan
matrix
\[
A_{\mathrm{II}}
= \begin{pmatrix} 2 & -2 & -2 \\ -2 & 2 & -2 \\ -2 & -2 & 2 \end{pmatrix},
\]
this condition is incompatible with the nonzero second real-root strings
$2\delta_i + \delta_j$, $i \ne j$.  The strict fixed-lift raw
$\Pi_X$-descent therefore cannot carry these brackets; they require the
normal-ordered extension $\widehat\Gamma_X$ and the additive homomorphism
$\overline\Pi_X$.
(proved, \cite{GN}.)
\end{theorem}

\begin{proof}
If the raw pushforward is $\Gamma_{\mathrm{gram}}$-graded, a nonzero
bracket of classes of raw Gram degrees $\Pi_X(c)$ and $\Pi_X(c')$ must
land in raw Gram degree $\Pi_X(c) + \Pi_X(c')$.  But the upstairs bracket
lands in physical degree $c + c'$, whose raw Gram shadow is
\[
\Pi_X(c + c') = \Pi_X(c) + \Pi_X(c') + B(c, c').
\]
Hence $B(c, c') = 0$ on every nonzero bracket channel.

Take physical lifts $c_i, c_j$ of distinct type-II real simple root
degrees $\gamma_i, \gamma_j$, $\Pi_X(c_i) = \gamma_i$ with $\gamma_i\ne
0$.  The Gritsenko--Nikulin Cartan matrix $A_{\mathrm{II}}$ has
off-diagonal entries $-2$, so distinct simple roots satisfy
$\langle\delta_i, \delta_j\rangle = -2$, and the real-root string
through $\delta_j$ in the $\delta_i$-direction includes the levels
$\delta_i + \delta_j$ and $2\delta_i + \delta_j$ (proved,
\cite[\S2--3]{GN}).  Both correspond to nonzero brackets, hence raw
pushforward forces
\[
B(c_i, c_j) = 0,\qquad
B(c_i, c_i + c_j) = 0.
\]
Bilinearity and polarization give
\[
B(c_i, c_i + c_j) = B(c_i, c_i) + B(c_i, c_j)
= 2\,\Pi_X(c_i) + 0 = 2\,\gamma_i,
\]
which is nonzero in the torsion-free lattice $\Gamma_{\mathrm{gram}}$.
Contradiction.
\end{proof}

\begin{remark}[Scope of the no-go]
\label{rem:no-go-scope}
Theorem~\ref{thm:raw-gram-descent-no-go} excludes the strict
\emph{homogeneous fixed-lift} raw pushforward.  It does not exclude
isolated $B$-isotropic bracket channels and does not decide the
mixed-lift fibre-summed construction in which several lifts of a single
Gram degree are summed before bracketing; the latter remains an open
mixed-lift question.  The normal-ordered extension removes exactly the
displayed obstruction because $\overline\Pi_X$ is additive on
$\widehat\Gamma_X$.  Chain-level Hall strictification
($A_\infty$-compatibility, cyclic equivariance, Frobenius,
multiplicative orientation) is a separate input system, named in
\S\ref{sec:k3xe-four-obstructions} and developed in
Chapter~\ref{chap:dirac-igusa-pro-object}.
\end{remark}

\begin{definition}[Gram fibres, zero fibre, and fibre-summed parity blocks]
\label{def:gram-fibres-zero-fibre}
For \(\gamma\in\Gamma_{\mathrm{gram}}\), set
\[
\Gamma_X(\gamma):=\Pi_X^{-1}(\gamma)\subset\Gamma_X^{\mathrm{form}}.
\]
The \emph{zero Gram fibre} is
\[
K_{\Pi,0}:=\Gamma_X(0,0,0)
=\{(Q,P)\mid Q^2=P^2=Q\cdot P=0\}.
\]
It is a fibre of the quadratic map \(\Pi_X\), not the kernel of a group
homomorphism.  If \(V=\bigoplus_{c\in A}V_c\) is a finite
\(\Z/2\)-graded formal source over a finite subset
\(A\subset\Gamma_X^{\mathrm{form}}\), the fibre-summed Gram pushforward
is the \(\Z/2\)-graded vector space
\[
(\Pi_{X,*}^{\mathrm{fib}}V)_{\gamma,\bar p}
:=\bigoplus_{\substack{c\in A\cap\Gamma_X(\gamma)}}V_{c,\bar p},
\qquad \bar p\in\Z/2.
\]
The signed superdimension
\(\dim(\Pi_{X,*}^{\mathrm{fib}}V)_{\gamma,\bar0}
-\dim(\Pi_{X,*}^{\mathrm{fib}}V)_{\gamma,\bar1}\)
is a scalar shadow of this two-integer datum.
\end{definition}

\begin{proposition}[The zero Gram fibre is not an additive subgroup]
\label{prop:zero-gram-fibre-not-subgroup}
\textnormal{[proved]}
In the two-hyperbolic-plane chart of
Lemma~\ref{lem:primitive-lift-every-gram-triple}, the zero Gram fibre
\(K_{\Pi,0}\) is not closed under addition.  Hence no quotient of
\(\Gamma_X^{\mathrm{form}}\) by \(K_{\Pi,0}\) is defined at the level of
abelian groups.
\end{proposition}

\begin{proof}
Let
\[
k_1=(e_1,e_2),\qquad k_2=(f_1,-f_2).
\]
Then
\[
\Pi_X(k_1)=\Pi_X(k_2)=(0,0,0),
\]
because \(e_1,e_2,f_1,f_2\) are isotropic and the two hyperbolic planes
are orthogonal.  Their sum is
\[
k_1+k_2=(e_1+f_1,e_2-f_2).
\]
The first component has square \(2\), the second has square \(-2\), and
the two components are orthogonal.  Therefore
\[
\Pi_X(k_1+k_2)=(1,0,-1)\ne(0,0,0).
\]
Thus \(K_{\Pi,0}\) is a null-cone fibre of a quadratic map, not a
subgroup.
\end{proof}

\begin{remark}[Fibre summation does not prove a Hall radical]
\label{rem:fibre-summation-not-radical}
Definition~\ref{def:gram-fibres-zero-fibre} proves only the formal
bookkeeping required before a finite Hall window can be compared with a
target root window.  A fibre with one even and one odd source vector has
signed superdimension \(0\) and parity profile \((1,1)\); it is not the
zero source degree.  The radical along \(K_{\Pi,0}\) is therefore a
separate Hopf-radical datum inside a finite source: it must specify the
pairing matrix, kernel matrix, quotient splitting, ideal/coideal
identities, and transition exactness.  It is not obtained by quotienting
the formal charge lattice by \(K_{\Pi,0}\).
\end{remark}

The packet \(\texttt{certificates/charge/gram\_fibre\_kernel}\) checks
these formal assertions on finite rows.  It records two distinct
formal charges with Gram label \((1,1,0)\), two with Gram label
\((1,1,1)\), and two nonzero charges in \(K_{\Pi,0}\) whose parity
profile is \((1,1)\) although the signed superdimension is \(0\).  It
also records the witness
\[
(e_1,e_2)+(f_1,-f_2)=(e_1+f_1,e_2-f_2),\qquad
\Pi_X(e_1+f_1,e_2-f_2)=(1,0,-1),
\]
so the zero fibre is not additive.  The packet does not prove finite
HN boundedness, compact Hall support, source Hall brackets, source
parity, a Hopf radical, exactness of a finite pushforward, Pfaffian
orientations, or protected traces.

\begin{definition}[Zero-fibre radical on a finite formal source]
\label{def:zero-fibre-radical-finite-source}
Let \(V=\bigoplus_{c\in A}V_c\) be a finite
\(\Gamma_X^{\mathrm{form}}\)-graded \(\Q\)-vector space with
\(\Z/2\)-parity, where \(A\subset\Gamma_X^{\mathrm{form}}\) is finite.
Let
\[
V_{\Pi,0}:=\bigoplus_{c\in A\cap K_{\Pi,0}}V_c.
\]
For a bilinear pairing \(h:V_{\Pi,0}\otimes V_{\Pi,0}\to\Q\), the
\emph{zero-fibre radical} is the ordinary linear-algebra kernel
\[
\operatorname{Rad}_{\Pi,0}(V,h)
:=\ker\!\left(V_{\Pi,0}\longrightarrow V_{\Pi,0}^{\vee}\right),
\qquad
v\longmapsto h(v,-).
\]
It is a subspace of the finite source vector space, not a quotient of
\(\Gamma_X^{\mathrm{form}}\) by the zero Gram fibre.
\end{definition}

\begin{proposition}[Exactness of finite fibre-summed Gram pushforward]
\label{prop:finite-fibre-pushforward-exact}
\textnormal{[proved]}
The functor
\[
\Pi_{X,*}^{\mathrm{fib}}:
\mathrm{Vect}^{\mathrm{fin}}_{\Gamma_X^{\mathrm{form}},\Z/2}
\longrightarrow
\mathrm{Vect}^{\mathrm{fin}}_{\Gamma_{\mathrm{gram}},\Z/2},
\qquad
(\Pi_{X,*}^{\mathrm{fib}}V)_{\gamma,\bar p}
=\bigoplus_{\Pi_X(c)=\gamma}V_{c,\bar p},
\]
is exact on finite-support graded \(\Q\)-vector spaces.  Consequently a
degreewise exact sequence
\[
0\longrightarrow A\longrightarrow B\longrightarrow C\longrightarrow0
\]
of finite formal sources remains exact after applying
\(\Pi_{X,*}^{\mathrm{fib}}\).
\end{proposition}

\begin{proof}
Fix \((\gamma,\bar p)\).  The \((\gamma,\bar p)\)-component of
\(\Pi_{X,*}^{\mathrm{fib}}\) is the finite direct sum of the
\((c,\bar p)\)-components over the fibre \(\Pi_X^{-1}(\gamma)\cap A\).
Finite direct sums are exact in the category of \(\Q\)-vector spaces.
Taking this finite direct sum in every Gram/parity component gives the
displayed exact sequence after pushforward.
\end{proof}

The packet
\(\texttt{certificates/charge/fibre\_pushforward\_exactness}\) checks
this finite linear algebra on six displayed Gram/parity blocks.  It
records matrices for \(0\to A\to B\to C\to0\), verifies that the
projection kernel equals the injection image in each block, verifies
the same equality after fibre-summed pushforward, and records a
zero-fibre pairing matrix whose radical is one-dimensional.  The packet
does not prove finite HN boundedness, compact Hall support, source Hall
brackets, Hopf ideal/coideal closure of the radical, exact Hall
pushforward, bar-construction commutation, Pfaffian orientations, or
protected traces.

\begin{proposition}[Finite normal-ordered pushforward preserves parity]
\label{prop:finite-parity-pushforward}
\textnormal{[proved]}
Let
\[
V=\bigoplus_{\hat\gamma\in A}
\left(V_{\hat\gamma,\bar0}\oplus V_{\hat\gamma,\bar1}\right)
\]
be a finite-support \(\widehat\Gamma_X\)-graded super vector space, and
let
\(\overline\Pi_X:\widehat\Gamma_X\to\Gamma_{\mathrm{gram}}\) be the
additive normal-ordered Gram map.  Define the pushed super vector space
by
\[
(\overline\Pi_{X,*}^{\mathrm{fib}}V)_{\gamma,\bar p}
:=
\bigoplus_{\substack{\hat\gamma\in A\\
\overline\Pi_X(\hat\gamma)=\gamma}}
V_{\hat\gamma,\bar p},
\qquad \bar p\in\Z/2.
\]
Then \(\overline\Pi_{X,*}^{\mathrm{fib}}\) preserves the even and odd
summands separately.  In particular the pair
\[
\bigl(\dim(\overline\Pi_{X,*}^{\mathrm{fib}}V)_{\gamma,\bar0},
\dim(\overline\Pi_{X,*}^{\mathrm{fib}}V)_{\gamma,\bar1}\bigr)
\]
is the fibrewise sum of the source parity dimensions, while the signed
superdimension is only its scalar shadow.
\end{proposition}

\begin{proof}
For fixed \((\gamma,\bar p)\), the displayed definition is an ordinary
finite direct sum of precisely the \(\bar p\)-parity components whose
normal-ordered Gram degree is \(\gamma\).  No even summand is placed in
an odd summand and no odd summand is placed in an even summand.  Thus
the pushed even block is the direct sum of source even blocks over the
fibre of \(\overline\Pi_X\), and the pushed odd block is the direct sum
of source odd blocks over the same fibre.  Taking the difference of the
two dimensions is a later scalar trace operation and does not recover
the two parity ranks.
\end{proof}

The packet \(\texttt{certificates/charge/parity\_pushforward}\) checks
this finite statement on three displayed Gram degrees.  It records the
two parity dimensions before and after pushforward and includes a
zero-degree row with parity profile \((1,1)\) and signed
superdimension \(0\).  It proves preservation of finite parity blocks
under the additive normal-ordered pushforward.  It does not prove the
source parity theorem, the Gritsenko--Nikulin/Kac target parity
equality, compact source representatives, primitive recognition,
Pfaffian orientations, or protected traces.

\begin{proposition}[Finite bar construction commutes with the
normal-ordered Gram pushforward]
\label{prop:finite-bar-pushforward-commutes}
\textnormal{[proved]}
Let \(A=\Q\mathbf 1\oplus\bar A\) be a finite-support augmented
associative algebra graded by the additive group
\(\widehat\Gamma_X\), with multiplication preserving the
\(\widehat\Gamma_X\)-degree.  Let
\[
T^c_{\le N}(\bar A)=
\bigoplus_{1\le p\le N}\bar A^{\otimes p}
\]
be the finite reduced tensor coalgebra with deconcatenation coproduct,
graded by the sum of the \(\widehat\Gamma_X\)-degrees of the letters.
Then the additive homomorphism
\(\overline\Pi_X:\widehat\Gamma_X\to\Gamma_{\mathrm{gram}}\) induces a
canonical isomorphism of finite
\(\Gamma_{\mathrm{gram}}\)-graded coalgebras
\[
\overline\Pi_{X,*}^{\mathrm{fib}}\,T^c_{\le N}(\bar A)
\;\cong\;
T^c_{\le N}\!\left(\overline\Pi_{X,*}^{\mathrm{fib}}\bar A\right).
\]
If the bar differential is defined by multiplication in \(A\), the
same isomorphism intertwines the finite bar differential.
\end{proposition}

\begin{proof}
For a word \([a_1|\cdots|a_p]\), the source degree is
\(\hat\gamma_1+\cdots+\hat\gamma_p\) in the additive group
\(\widehat\Gamma_X\).  Since \(\overline\Pi_X\) is a group homomorphism,
\[
\overline\Pi_X(\hat\gamma_1+\cdots+\hat\gamma_p)
=\overline\Pi_X(\hat\gamma_1)+\cdots+\overline\Pi_X(\hat\gamma_p).
\]
Thus pushing the word degree forward after forming the tensor coalgebra
and forming the tensor coalgebra after pushing the letters forward give
the same finite direct sum in each
\(\Gamma_{\mathrm{gram}}\)-degree.  Deconcatenation cuts a word into
two subwords; the displayed equality is additive on the two subword
degrees, so the coproduct is preserved.  The internal bar differential
replaces adjacent letters \(a_i,a_{i+1}\) by their product.  Degree
preservation of multiplication in \(A\) gives
\(\deg(a_i a_{i+1})=\deg(a_i)+\deg(a_{i+1})\), hence the same
\(\Gamma_{\mathrm{gram}}\)-degree after \(\overline\Pi_X\).  Therefore
the finite bar differential is intertwined.
\end{proof}

The packet
\(\texttt{certificates/charge/bar\_pushforward\_commutation}\) checks
this statement on finite words of length at most three.  It records
letter degrees, word degrees, deconcatenation splits, and product terms
in the bar differential, all using the additive normal-ordered Gram
grading.  It does not certify the raw quadratic \(\Pi_X\)-pushforward,
the chiral Koszul quasi-isomorphism, Hall bracket compatibility, Hopf
pairing compatibility, Pfaffian orientations, or protected traces.

\begin{proposition}[Finite bracket and pairing degrees descend under
normal-ordered Gram pushforward]
\label{prop:finite-bracket-pairing-pushforward}
\textnormal{[proved]}
Let \(P=\bigoplus_{\hat\gamma\in\widehat\Gamma_X}P_{\hat\gamma}\) be a
finite-support \(\widehat\Gamma_X\)-graded Lie super vector space with
\[
[P_{\hat\gamma},P_{\hat\delta}]
\subset P_{\hat\gamma+\hat\delta}.
\]
Then \(\overline\Pi_{X,*}^{\mathrm{fib}}P\) is a
\(\Gamma_{\mathrm{gram}}\)-graded Lie super vector space.  If
\[
h:P\otimes P\to\Q
\]
is a homogeneous degree-zero bilinear form, in the sense that
\[
h(P_{\hat\gamma},P_{\hat\delta})=0
\quad\text{unless}\quad \hat\gamma+\hat\delta=0,
\]
then the fibre-summed form on
\(\overline\Pi_{X,*}^{\mathrm{fib}}P\) is homogeneous of degree zero
for the \(\Gamma_{\mathrm{gram}}\)-grading.
\end{proposition}

\begin{proof}
For \(x\in P_{\hat\gamma}\) and \(y\in P_{\hat\delta}\), the source
bracket has degree \(\hat\gamma+\hat\delta\).  Applying the additive
homomorphism \(\overline\Pi_X\) gives
\[
\overline\Pi_X(\hat\gamma+\hat\delta)
=\overline\Pi_X(\hat\gamma)+\overline\Pi_X(\hat\delta),
\]
so the same bracket is graded after fibre-summed pushforward.  The
Jacobi identity and super-skew symmetry are unchanged, since the
underlying bracket is the same linear map on the same finite direct
sum.

If \(h(x,y)\ne0\), then \(\hat\gamma+\hat\delta=0\).  Applying
\(\overline\Pi_X\) gives
\(\overline\Pi_X(\hat\gamma)+\overline\Pi_X(\hat\delta)=0\).  Thus the
pushed pairing can be nonzero only on opposite Gram degrees.  This is
degree-zero homogeneity of the pushed form.  No Hopf adjointness,
Frobenius cyclic identity, quotient non-degeneracy, or primitive
closure is used.
\end{proof}

The packet
\(\texttt{certificates/charge/hall\_pairing\_pushforward\_compatibility}\)
checks this finite algebraic statement on supplied bracket and pairing
rows.  It verifies that three bracket rows have additive
normal-ordered Gram degree and that three nonzero pairing rows have
total Gram degree \(0\).  It does not construct compact Hall
correspondences, prove primitive closure, prove Hopf adjointness, prove
the Frobenius cyclic identity, prove quotient non-degeneracy, prove
Pfaffian orientations, or prove protected traces.

\begin{proposition}[Finite normal-ordered pushforward preserves the PBW filtration]
\label{prop:finite-pbw-pushforward}
\textnormal{[proved]}
Let \(U=\bigoplus_{\hat\gamma\in A}U_{\hat\gamma}\) be a
finite-support \(\widehat\Gamma_X\)-graded vector space equipped with
an increasing PBW filtration \(F_{\le d}U\) by homogeneous subspaces:
\[
F_{\le d}U
=\bigoplus_{\hat\gamma\in A}
\bigl(F_{\le d}U\cap U_{\hat\gamma}\bigr).
\]
Write \(F_{\le d}U_{\hat\gamma}:=F_{\le d}U\cap U_{\hat\gamma}\).
Define
\[
F_{\le d}
(\overline\Pi_{X,*}^{\mathrm{fib}}U)_{\gamma}
:=
\bigoplus_{\substack{\hat\gamma\in A\\
\overline\Pi_X(\hat\gamma)=\gamma}}
F_{\le d}U_{\hat\gamma}.
\]
Then \(\overline\Pi_{X,*}^{\mathrm{fib}}\) preserves the PBW
filtration, and for every \(d\)
\[
\operatorname{gr}^{\mathrm{PBW}}_d
(\overline\Pi_{X,*}^{\mathrm{fib}}U)_{\gamma}
\cong
\bigoplus_{\substack{\hat\gamma\in A\\
\overline\Pi_X(\hat\gamma)=\gamma}}
\operatorname{gr}^{\mathrm{PBW}}_d U_{\hat\gamma}.
\]
If \(U\) is a filtered algebra whose multiplication has
\(\widehat\Gamma_X\)-degree \(\hat\gamma+\hat\delta\) and PBW degree
\(\le a+b\) on \(F_{\le a}U_{\hat\gamma}\cdot
F_{\le b}U_{\hat\delta}\), the pushed multiplication has
\(\Gamma_{\mathrm{gram}}\)-degree
\(\overline\Pi_X(\hat\gamma)+\overline\Pi_X(\hat\delta)\) and PBW
degree \(\le a+b\).
\end{proposition}

\begin{proof}
For fixed \((\gamma,d)\), the displayed formula is a finite direct sum
over the fibre of \(\overline\Pi_X\).  Since
\(F_{\le d}U_{\hat\gamma}\subset F_{\le d+1}U_{\hat\gamma}\) for each
\(\hat\gamma\), the same inclusion holds after taking the finite direct
sum over the fibre.  Finite direct sums are exact over \(\Q\), hence
the quotient \(F_{\le d}/F_{\le d-1}\) commutes with the displayed
pushforward and gives the associated-graded isomorphism.  For the
algebra statement, the source multiplication sends
\((\hat\gamma,\hat\delta)\) to \(\hat\gamma+\hat\delta\) and has PBW
degree \(\le a+b\); applying the additive homomorphism
\(\overline\Pi_X\) sends the Gram degree to
\(\overline\Pi_X(\hat\gamma)+\overline\Pi_X(\hat\delta)\) and does not
change the filtration index.  Thus the pushed filtered multiplication
has the asserted degree and filtration bound.
\end{proof}

The packet \(\texttt{certificates/charge/pbw\_pushforward}\) checks
this finite filtered statement on three displayed Gram degrees.  It
records cumulative PBW dimensions, verifies monotonicity of the
filtration, and recomputes the associated-graded dimensions as
successive quotients.  It proves only preservation of a supplied finite
PBW filtration under the additive normal-ordered pushforward.  It does
not prove the compact-source PBW comparison, target PBW comparison,
no-extra-relations theorem, primitive recognition, Pfaffian
orientations, or protected traces.

\begin{proposition}[Finite normal-ordered pushforward preserves triangular decomposition]
\label{prop:finite-triangular-pushforward}
\textnormal{[proved]}
Let \(h:\Gamma_{\mathrm{gram}}\to\Z\) be an additive height functional.
For a finite-support \(\widehat\Gamma_X\)-graded vector space
\(V=\bigoplus_{\hat\gamma\in A}V_{\hat\gamma}\), define
\[
V^+=\bigoplus_{h(\overline\Pi_X(\hat\gamma))>0}V_{\hat\gamma},
\qquad
V^0=\bigoplus_{h(\overline\Pi_X(\hat\gamma))=0}V_{\hat\gamma},
\qquad
V^-=\bigoplus_{h(\overline\Pi_X(\hat\gamma))<0}V_{\hat\gamma}.
\]
Then
\[
\overline\Pi_{X,*}^{\mathrm{fib}}V
=
(\overline\Pi_{X,*}^{\mathrm{fib}}V)^-
\oplus
(\overline\Pi_{X,*}^{\mathrm{fib}}V)^0
\oplus
(\overline\Pi_{X,*}^{\mathrm{fib}}V)^+
\]
with each summand obtained by pushing forward the corresponding source
summand over the same fibres of \(\overline\Pi_X\).  If \(V\) carries a
homogeneous bracket
\[
[V_{\hat\gamma},V_{\hat\delta}]
\subset V_{\hat\gamma+\hat\delta},
\]
then the pushed bracket has height
\[
h\!\left(\overline\Pi_X(\hat\gamma+\hat\delta)\right)
=h(\overline\Pi_X(\hat\gamma))
+h(\overline\Pi_X(\hat\delta)).
\]
Thus positive-positive brackets remain positive, negative-negative
brackets remain negative, Cartan-positive and Cartan-negative brackets
remain in the corresponding sign, and positive-negative brackets land
in the sign determined by the actual height sum.
\end{proposition}

\begin{proof}
The three displayed subspaces form a direct sum because every integer
height is exactly one of positive, zero, or negative.  For a fixed
Gram degree \(\gamma\), the fibre-summed pushforward takes the finite
direct sum over \(\overline\Pi_X^{-1}(\gamma)\).  All source summands
in this fibre have the same value \(h(\gamma)\), hence the same
triangular sign.  Therefore no positive source summand is pushed into
the Cartan or negative part, and similarly for the other two signs.

For the bracket statement, source homogeneity gives degree
\(\hat\gamma+\hat\delta\).  Additivity of \(\overline\Pi_X\) and of
\(h\) gives the displayed height equality.  The triangular sign after
pushforward is therefore the sign of this integer.  The four stated
same-sign and Cartan-action cases are immediate; the
positive-negative case is not collapsed to a universal rule, but is
read from the actual height sum.
\end{proof}

The packet
\(\texttt{certificates/charge/triangular\_pushforward}\) checks this
finite statement on seven displayed source blocks and five supplied
bracket sign rows.  It verifies that the supplied integer height agrees
with the positive, Cartan, and negative labels, that the fibre-summed
pushforward keeps the three direct-sum parts separate, and that the
displayed bracket rows land in the triangular part determined by the
pushed height.  It does not prove exact triangular completion, target
triangular completion, compact-source primitive recognition,
no-extra-relations, PBW comparison, Pfaffian orientations, or protected
traces.

\begin{remark}[Quadratic obstruction; no linear trivialisation]
\label{rem:quadratic-obstruction}
Suppose $B = \delta L$ for some additive linear $L : \Gamma_X^{\mathrm
{form}}\to\Gamma_{\mathrm{gram}}$.  Setting $c' = c$, $B(c, c) = 2 L(c)
- L(2 c)$; additivity gives $L(2 c) = 2 L(c)$, so $B(c, c) = 0$,
contradicting $B(c, c) = 2\,\Pi_X(c) \ne 0$ for any $c$ with
$(Q, Q), (Q, P), (P, P)$ not all zero.  Hence no additive linear
trivialisation exists.  The quadratic primitive $\Pi_X$ is necessary;
trivializing the additive cocycle requires a quadratic cochain, not a
linear one (proved).
\end{remark}

\subsection{Normal-ordering flag formula}
\label{subsec:normal-ordering-flag-formula}

\begin{lemma}[Multiplication of raw placements]
\label{lem:multiplication-of-raw-placements}
For $c_1, \ldots, c_k \in \Gamma_X^{\mathrm{form}}$,
\[
i_0(c_1)\star\cdots\star i_0(c_k)
= \biggl(\sum_i c_i,\ \sum_{i < j} B(c_i, c_j)\biggr),
\]
and consequently
\[
\overline\Pi_X\!\left(i_0(c_1)\star\cdots\star i_0(c_k)\right)
= \sum_i \Pi_X(c_i).
\]
(proved.)
\end{lemma}

\begin{proof}
Induct on $k$.  The case $k = 2$ is the definition of $\star$.  At
step $k\to k + 1$,
\[
B\!\left(\sum_{i\le k} c_i,\,c_{k + 1}\right)
= \sum_{i\le k} B(c_i,\,c_{k + 1})
\]
by bilinearity.  Applying $\overline\Pi_X$ subtracts the accumulated
pairwise $B$-term from $\Pi_X(\sum_i c_i) = \sum_i \Pi_X(c_i) +
\sum_{i < j} B(c_i, c_j)$.
\end{proof}

\begin{remark}[Normal-ordered primitive recognition target]
\label{rem:normal-ordered-recognition-target}
The interface with the BKM denominator algebra is the normal-ordered
charge route
\[
\Gamma_X^{\mathrm{form}}\hookrightarrow\widehat\Gamma_X
\xrightarrow{\overline\Pi_X}\Gamma_{\mathrm{gram}}
\xrightarrow{\alpha}\Lambda_{II}^{2,1},
\]
where $\alpha(n, l, m) = 2 n f_2 - l f_3 + 2 m f_{-2}$ is the type-IV
embedding into the BKM root lattice (proved, \cite{GN}).  At
chain level this becomes the descent $\overline\Pi_{X, *}^{\Theta}$ of
Chapter~\ref{chap:dirac-igusa-pro-object}.  The primitive recognition
problem is to construct a protected primitive Hall algebra whose
normal-ordered descended bracket is the Borcherds--Kac bracket of
$\mathfrak g_{\Delta_5}$.  Raw Gram pushforward is not an admissible
recognition target.
\end{remark}

\section{The hybrid $\operatorname{Ran}^{\mathrm{hyb}}(E)$ factorization
carrier}
\label{sec:k3xe-hybrid-ran-carrier}

After integration along the K3 direction, the chiral object is supported
on the elliptic factor $E$.  The ordinary finite-support carrier is the Ran space
$\operatorname{Ran}(E)$ of finite subsets, but a one-dimensional
effective cycle with positive elliptic degree projects onto all of $E$.  The
hybrid carrier $\operatorname{Ran}^{\mathrm{hyb}}(E)$ is a finite-stage
prestack whose \(b=0\) local and \(b>0\) wrapped subprestacks are
separated before the mixed local/wrapped correspondences are supplied as
additional finite-stage data.
At finite stage those data are represented by the tables
\(\mathrm{hybrid\_ran\_prestack\_definition}\),
\(\mathrm{local\_stratum\_b0\_prestack\_definition}\),
\(\mathrm{wrapped\_stratum\_bpositive\_prestack\_definition}\),
\(\mathrm{geometric\_elliptic\_degree\_map}\),
\(\mathrm{geometric\_elliptic\_degree\_additivity}\),
\(\mathrm{geometric\_borcherds\_degree\_separation}\),
\(\mathrm{positive\_elliptic\_degree\_projection}\),
\(\mathrm{hybrid\_base}\), \(\mathrm{local\_configurations}\),
\(\mathrm{wrapped\_prequotients}\),
\(\mathrm{local\_local\_correspondence\_definition}\),
\(\mathrm{mixed\_local\_wrapped\_correspondence\_definition}\),
\(\mathrm{wrapped\_wrapped\_correspondence\_definition}\),
\(\mathrm{two\_sided\_mixed\_correspondence\_definition}\),
\(\mathrm{correspondences}\),
\(\mathrm{flag\_atlas}\), \(\mathrm{higher\_coloured}\),
\(\mathrm{quotient\_pseudofunctor}\), \(\mathrm{protected\_integration}\),
\(\mathrm{transitions}\), and \(\mathrm{scalar\_firewall}\).  They are
not supplied by the ordinary Ran space, a quotient-first Hall product,
or the Borcherds \(s\)-degree.

\subsection{Projection-locality obstruction}
\label{subsec:projection-locality-obstruction}

\begin{lemma}[Projection-locality obstruction]
\label{lem:projection-locality-obstruction}
Let
\[
Z=\sum_i m_i[C_i]\in Z_1(S\times E),\qquad m_i>0,
\]
be an effective one-cycle.  If \((p_E)_*Z=b[E]\) with \(b>0\), then
\(p_E(|Z|)=E\).  Hence \(|Z|\) is not local on
\(\operatorname{Ran}(E)\).
\end{lemma}

\begin{proof}
For each irreducible component \(C_i\), the proper image \(p_E(C_i)\)
is either a point or a closed one-dimensional subset of the irreducible
curve \(E\).  In the second case it is all of \(E\), and
\((p_E)_*[C_i]=\deg(C_i/E)[E]\).  Since the effective sum
\((p_E)_*Z=b[E]\) has \(b>0\), at least one component maps dominantly
to \(E\).  Therefore \(p_E(|Z|)=E\).
\end{proof}

\begin{remark}[Geometric versus Borcherds elliptic degree]
\label{rem:geometric-vs-borcherds-degree}
The lemma fixes the necessity of a wrapped stratum.  The geometric
elliptic-degree map $b_R^{\mathrm{geom}}$, on a finite Hall stage $R$,
records the integer $\deg_E$ of a chosen retained class.  The trace
exponent $m_R^{\mathrm{Bch}} := \mathrm{pr}_3\overline\Pi_X$, on the same stage,
records the third coordinate of the normal-ordered Gram pushforward.
The two are not equal as functions on the formal lattice; the Borcherds
exponent $m$ records the magnetic Mukai pairing $(P, P)/2$, which sees
$\Lambda_S$-square data, while $b_R^{\mathrm{geom}}$ records the
$E$-degree of the geometric support, which sees $H^0(S)\oplus H^4(S)$
data through $r$ and $\mathrm{ch}_2$ of the K3 factor.  The wrapped
stratum carries $b_R^{\mathrm{geom}} > 0$ branes; the local stratum
carries $b_R^{\mathrm{geom}} = 0$ branes whose magnetic component lies
in the projection-finite kernel.  This is the
$\operatorname{Ran}(E)$-support-locality obstruction in the
K3-integrated chiral shadow.
\end{remark}

\subsection{The hybrid carrier}
\label{subsec:hybrid-carrier-definition}

\begin{definition}[Hybrid Ran carrier]
\label{def:hybrid-ran-carrier}
At finite HN height \(R\), the \emph{hybrid Ran carrier} is the
prestack
\[
\operatorname{Ran}^{\mathrm{hyb}}_{R,\mathrm{pre}}(E):
\mathrm{Aff}_{\C}^{op}\to\infty\textup{-}\mathrm{Gpd}
\]
of Definition~\ref{def:hybrid-ran-prestack}.  A test family consists
of finitely many local maps to \(E\), retained local colour labels,
finitely many wrapped colour labels, and \(T\)-families in the
corresponding rigidified wrapped prequotient stacks with anchor maps to
\(E\).  Its local \(b=0\) stratum is the subprestack
\(\operatorname{Ran}^{\mathrm{loc}}_{R,\mathrm{pre}}(E)\) of
Definition~\ref{def:hybrid-local-stratum-prestack}.  Its wrapped
\(b>0\) stratum is the subprestack
\(\operatorname{Ran}^{\mathrm{wr}}_{R,\mathrm{pre}}(E)\) of
Definition~\ref{def:hybrid-wrapped-stratum-prestack}.  The stack
\(\operatorname{Ran}^{\mathrm{hyb}}_{R,\tau}(E)\) exists only after
descent rows for the chosen topology \(\tau\) have been supplied.  The
set
\[
\bigsqcup_{b\ge0}\operatorname{Ran}(E)\times\{b\}
\]
is the degree shadow, not the carrier.
\end{definition}

\begin{remark}[Two strata, two registers]
\label{rem:two-strata-two-registers}
On the $b = 0$ stratum the standard Ran/factorization formalism applies
and supplies the local chiral factorization input in
$\mathrm{Fact}(\operatorname{Ran}(E))$ (folklore,
\cite{BD,FrancisGaitsgoryCKD}).  On the $b > 0$
stratum the support of $\mathrm{Supp}\,\mathcal F$ projects onto $E$, so
the brane is global; its data live on the wrapped prequotient and after
$E$-quotient become $E$-equivariant data on the elliptic curve.  The
hybrid carrier is the smallest carrier reflecting both strata; it is
not a single Ran space.
\end{remark}

\subsection{Mixed local/wrapped correspondences}
\label{subsec:mixed-correspondences}

\begin{definition}[Mixed local/wrapped correspondences (finite Hall
stage $R$)]
\label{def:mixed-correspondences}
Fix a finite Hall stage $R$ and retained Liu--Hilbert classes $\Xi_A,
\Xi_B, \Xi_C$ with $\Xi_C = \Xi_A + \Xi_B$ in the closure rules of
\cite{KSCoHA}.  The \emph{retained extension correspondence}
\[
\mathfrak E^{\mathrm{ret}}_{\Xi_A, \Xi_B; \Xi_C}
\xrightarrow{p}
\mathfrak M^{\mathrm{ss}}_{\Xi_A}\times\mathfrak M^{\mathrm{ss}}_{\Xi_B},
\qquad
\mathfrak E^{\mathrm{ret}}_{\Xi_A, \Xi_B; \Xi_C}
\xrightarrow{q}
\mathfrak M^{\mathrm{ss}}_{\Xi_C}
\]
is the closed derived flag-Quot/filtration stack over
$\mathfrak M^{\mathrm{ss}}_{\Xi_C}$, with $q$ required at this finite
stage to be proper or $\mathrm{Coh}^{\mathrm{prop}}$-admissible in the
compactified model (folklore).  Three colours of
correspondence appear,
indexed by the location of $\Xi_A, \Xi_B$ in
$\operatorname{Ran}^{\mathrm{hyb}}(E)$:
\begin{enumerate}[label = \textup{(LL)},leftmargin = *]
\item[\textup{(LL)}] \emph{local-local}: $b_A = b_B = 0$;
\item[\textup{(LW)}] \emph{ordered local/wrapped}: $b_A = 0$, $b_B > 0$;
\item[\textup{(WW)}] \emph{wrapped-wrapped}: $b_A, b_B > 0$.
\end{enumerate}
The local-local row is the ordered projection-finite diagram
\[
\mathfrak E^{\mathrm{LL}}_{\alpha,\beta;\gamma,R,I,J,K}
\longrightarrow
\mathfrak M^{\mathrm{loc,cl}}_{\alpha,R,I}
\times
\mathfrak M^{\mathrm{loc,cl}}_{\beta,R,J}
\longrightarrow
\mathfrak M^{\mathrm{loc,cl}}_{\gamma,R,K},
\]
where \(\alpha:I\to\Gamma_R^{\mathrm{loc}}\),
\(\beta:J\to\Gamma_R^{\mathrm{loc}}\),
\(\gamma=|\alpha|+|\beta|\) is retained in
\(\Gamma_R^{\mathrm{loc}}\), and \(K=I\sqcup J\) modulo Ran
collisions.  This is
Definition~\ref{def:local-local-correspondence-datum}; it has no
wrapped input and no reduced \(E\)-quotient.
The one-sided mixed row is the ordered pair of unreduced diagrams
\[
\mathfrak E^{\mathrm{LW}}_{\alpha,\eta;\zeta,R,I}
\longrightarrow
\mathfrak M^{\mathrm{loc,cl}}_{\alpha,R,I}
\times
\mathcal M_{\eta,R}^{\mathrm{wr,rig}}
\longrightarrow
\mathcal M_{\zeta,R}^{\mathrm{wr,rig}},
\]
and
\[
\mathfrak E^{\mathrm{WL}}_{\eta,\alpha;\zeta,R,I}
\longrightarrow
\mathcal M_{\eta,R}^{\mathrm{wr,rig}}
\times
\mathfrak M^{\mathrm{loc,cl}}_{\alpha,R,I}
\longrightarrow
\mathcal M_{\zeta,R}^{\mathrm{wr,rig}},
\]
where \(\alpha:I\to\Gamma_R^{\mathrm{loc}}\),
\(\eta\in\Gamma_R^{\mathrm{wr}}\), and
\(\zeta=|\alpha|+\eta\) or \(\eta+|\alpha|\) is retained in
\(\Gamma_R^{\mathrm{wr}}\).  This is
Definition~\ref{def:mixed-local-wrapped-correspondence-datum}; it is
not the local-local row, not the wrapped-wrapped row, not the two-sided
mixed row, and not a quotient-first product.
The wrapped-wrapped row is the ordered unreduced diagram
\[
\mathfrak E^{\mathrm{WW}}_{\eta_1,\eta_2;\zeta,R}
\longrightarrow
\mathcal M_{\eta_1,R}^{\mathrm{wr,rig}}
\times
\mathcal M_{\eta_2,R}^{\mathrm{wr,rig}}
\longrightarrow
\mathcal M_{\zeta,R}^{\mathrm{wr,rig}},
\]
where \(\eta_1,\eta_2\in\Gamma_R^{\mathrm{wr}}\) and
\(\zeta=\eta_1+\eta_2\) is retained in
\(\Gamma_R^{\mathrm{wr}}\).  This is
Definition~\ref{def:wrapped-wrapped-correspondence-datum}; it is
ordered, unreduced, and separate from symmetric descent and
source/target-anchor memory.
The ordered two-sided mixed row is the unreduced local--wrapped--local
diagram
\[
\mathfrak E^{\mathrm{LWL}}_{\alpha_-;\eta;\beta_+,R,I_-,I_+}
\longrightarrow
\mathfrak M^{\mathrm{loc,cl}}_{\alpha_-,R,I_-}
\times
\mathcal M_{\eta,R}^{\mathrm{wr,rig}}
\times
\mathfrak M^{\mathrm{loc,cl}}_{\beta_+,R,I_+}
\longrightarrow
\mathcal M_{\zeta,R}^{\mathrm{wr,rig}},
\]
where
\(\zeta=|\alpha_-|+\eta+|\beta_+|\in\Gamma_R^{\mathrm{wr}}\).
	This is
	Definition~\ref{def:two-sided-mixed-correspondence-datum}; it is not
	the full eight-word associativity atlas and not a quotient-first
	product.
	The source/target anchor-memory row is separate: for \(LW\), \(WL\),
	\(WW\), and \(LWL\) it records the ordered wrapped source anchors and
	wrapped target anchor on the unreduced correspondence stack, before the
	reduced \(E\)-quotient, as in
	Definition~\ref{def:source-target-anchor-memory-datum}.  The
	local/local row has no wrapped anchor memory.  This row does not
	construct the determinant anchor, compute its translation weight,
	repair the \(\chi=0\) degeneracy, or define \(o^\lambda_R\).
	The arity-three shadow is the eight-word two-step binary flag atlas
	\[
	\{L, W\}^3 = \{LLL, LLW, LWL, LWW, WLL, WLW, WWL, WWW\},
	\]
which records binary associativity only; higher-coloured tree
configurations live on a separate coloured tree category
$\mathrm{Tree}^{\mathrm{col}}_R$, supplied as additional input
(folklore, \cite{CYQGVolIII}).
\end{definition}

\begin{remark}[Eight binary words versus colored trees]
\label{rem:eight-words-vs-colored-trees}
The eight binary words of Definition~\ref{def:mixed-correspondences}
exhaust the arity-three two-step compositions $\{L, W\}^3$, but they do
not catalogue higher-coloured tree configurations, units, symmetry/order
conventions, refinement maps, descent, or overlaps.  Those data are
additional finite Hall input named at \S\ref{sec:k3xe-four-obstructions};
asserting them prematurely confuses binary associativity with the full
factorization-coherence structure (\cite{CYQGVolIII},
correspondence formalism).
\end{remark}

\subsection{Wrapped anchor and translation linearization}
\label{subsec:e-quotient-anchor}

\begin{definition}[Determinant anchor on the wrapped stratum]
\label{def:determinant-anchor}
Fix a finite HN height \(R\), a wrapped colour
\(\eta\in\Gamma_R^{\mathrm{wr}}\), and a test scheme \(T\).  Let
\(\mathcal F_{\eta,T}\) be the retained universal perfect complex on
\(X\times T\) pulled back from the wrapped prequotient row, and write
\[
\pi_{E,T}:X\times T\longrightarrow E\times T .
\]
Let \(\chi_\eta=\chi(\mathcal F_{\eta,t})\), constant on the retained
colour row.  The Knudsen--Mumford determinant functor
\cite{KnudsenMumfordDet,DeligneDetCohomology} gives a line bundle
\[
\mathscr D_{\eta,T}
:=\det R\pi_{E,T,*}\mathcal F_{\eta,T}
\]
on \(E\times T\), functorial in \(T\) as a point of the relative
Picard functor.  Since \(E\) has genus \(1\),
\(\deg\det C=\chi(C)\) for a perfect complex \(C\) on \(E\).  Hence
\(\mathscr D_{\eta,T}\) has relative degree \(\chi_\eta\), and
\[
\lambda^{\det}_{\eta,R}(\mathcal F_{\eta,T})
:=\mathscr D_{\eta,T}\otimes
\operatorname{pr}_E^*\mathcal O_E(-\chi_\eta\cdot0_E)
\]
where \(\operatorname{pr}_E:E\times T\to E\) is projection.  This line
bundle has relative degree \(0\).  Its class defines the
determinant anchor
\[
\lambda^{\det}_{\eta,R}:T\longrightarrow
\operatorname{Pic}^0(E)\simeq E .
\]
This constructs only the normalized determinant anchor.  The next
propositions compute its \(E\)-translation weight and isolate the
\(\chi_\eta=0\) degeneracy.  The finite-cover normalization to unit
weight, extra anchor data, anchor residual definition and vanishing,
quotient descent, and transition compatibility are separate rows.
\end{definition}

The packet
\(\texttt{certificates/hybrid/determinant\_anchor\_construction}\)
verifies
\(\texttt{DETERMINANT\_ANCHOR\_CONSTRUCTION\_VERIFIED}\): the
determinant-of-cohomology line is functorial in families, the
normalization has relative degree zero, and the resulting map lands in
\(\operatorname{Pic}^0(E)\simeq E\).  It does not compute the
translation weight, handle the \(\chi=0\) degeneracy, add extra anchor
data, define \(o^\lambda_R\), or prove quotient descent.

\begin{proposition}[Translation weight of the determinant anchor]
\label{prop:determinant-anchor-translation-weight}
\textnormal{[proved]}
Let \(a:T\to E\) be a section and let
\(\tau_a:E\times T\to E\times T\) be translation by \(a\).  Use the
support-translation convention
\[
a\cdot \mathcal F_{\eta,T}
:=(\operatorname{id}_S\times\tau_a)_*\mathcal F_{\eta,T}
\]
on \(X=S\times E\).  If \(\Gamma_a,\Gamma_0\subset E\times T\) are
the graphs of \(a\) and \(0_E\), then in the relative Picard functor
\[
\lambda^{\det}_{\eta,R}(a\cdot\mathcal F_{\eta,T})
\simeq
\lambda^{\det}_{\eta,R}(\mathcal F_{\eta,T})\otimes
\mathcal O_{E\times T}\!\bigl(\chi_\eta(\Gamma_a-\Gamma_0)\bigr).
\]
Under \(\operatorname{Pic}^0(E)\simeq E\), the induced translation law is
\[
\lambda^{\det}_{\eta,R}(a\cdot\mathcal F_{\eta,T})
=
\lambda^{\det}_{\eta,R}(\mathcal F_{\eta,T})+\chi_\eta a .
\]
Thus the \(E\)-translation weight of the determinant anchor is
\(\chi_\eta\).  Its kernel is \(E[\chi_\eta]\) when \(\chi_\eta\ne0\),
and all of \(E\) when \(\chi_\eta=0\).
\end{proposition}

\begin{proof}
Set
\(\mathscr D_{\eta,T}=\det R\pi_{E,T,*}\mathcal F_{\eta,T}\).  The
Cartesian square for translation on the elliptic factor gives, in the
relative Picard group,
\[
\det R\pi_{E,T,*}(a\cdot\mathcal F_{\eta,T})
\simeq
\tau_{a,*}\mathscr D_{\eta,T}
\simeq
\tau_{-a}^{*}\mathscr D_{\eta,T}.
\]
For a relative line bundle \(L\) of degree \(d\) on \(E\times T\),
translation acts trivially on the \(\operatorname{Pic}^0(E)\)-class and
shifts the degree part by
\[
\tau_{-a}^{*}L\otimes L^{-1}
\simeq
\mathcal O_{E\times T}\!\bigl(d(\Gamma_a-\Gamma_0)\bigr)
\quad
\text{in }\operatorname{Pic}^0(E\times T/T).
\]
Taking \(L=\mathscr D_{\eta,T}\) and \(d=\chi_\eta\), and then tensoring
both sides with the fixed normalization
\(\mathcal O_{E\times T}(-\chi_\eta\Gamma_0)\), gives the displayed
identity.  The sign is positive because the manuscript uses pushforward
by the support translation \(x\mapsto x+a\); the pullback convention
would reverse the sign.  The kernel statement is the kernel of the
multiplication homomorphism \([\chi_\eta]:E\to E\).
\end{proof}

The packet
\(\texttt{certificates/hybrid/determinant\_anchor\_translation\_weight}\)
verifies
\(\texttt{DETERMINANT\_ANCHOR\_TRANSLATION\_WEIGHT\_VERIFIED}\):
for the support-translation convention
\(a\cdot\mathcal F=(\operatorname{id}_S\times\tau_a)_*\mathcal F\),
the determinant anchor has translation weight \(\chi_\eta\), with
restriction \(h\mapsto\chi_\eta h\) on any finite stabilizer
\(H\subset E[N]\).  It does not choose a coherent stabilizer
linearization, make the weight one by finite cover, handle the
\(\chi_\eta=0\) degeneracy, add extra anchor data, define
\(o^\lambda_R\), prove quotient descent, or prove transition
compatibility.

\begin{proposition}[Zero-Euler determinant-anchor degeneracy]
\label{prop:determinant-anchor-chi-zero-degeneracy}
\textnormal{[proved]}
Assume \(\chi_\eta=0\).  Then for every section \(a:T\to E\),
\[
\lambda^{\det}_{\eta,R}(a\cdot\mathcal F_{\eta,T})
=
\lambda^{\det}_{\eta,R}(\mathcal F_{\eta,T})
\quad\text{in }\operatorname{Pic}^0(E).
\]
On a retained finite row where the \(E\)-translation action is free, a
geometric point \(a\ne0\) gives distinct objects
\(\mathcal F_{\eta,t}\) and \(a\cdot\mathcal F_{\eta,t}\) with the same
determinant anchor.  Hence the determinant anchor has full
\(E\)-orbit kernel on the \(\chi_\eta=0\) branch and has
one-dimensional separation defect.  A final wrapped anchor on this
branch must contain extra anchor data not determined by
\(\lambda^{\det}_{\eta,R}\).
\end{proposition}

\begin{proof}
Proposition~\ref{prop:determinant-anchor-translation-weight} gives
\[
\lambda^{\det}_{\eta,R}(a\cdot\mathcal F_{\eta,T})
=
\lambda^{\det}_{\eta,R}(\mathcal F_{\eta,T})+\chi_\eta a .
\]
When \(\chi_\eta=0\) this is the displayed equality for all
\(a\in E\).  The retained \(E\)-translation-rigidification datum is
free on the retained row by
Proposition~\ref{prop:retained-e-translation-rigidifications}; hence a
nonzero geometric translation gives a different point in the
prequotient before the reduced quotient is taken.  The two points have
the same determinant anchor.  Therefore determinant-anchor-only data
does not separate the \(E\)-orbit.  The orbit is one-dimensional, so
the class-level separation defect has rank \(1\).  Extra anchor data is
therefore a necessary input for the \(\chi_\eta=0\) branch, not a
consequence of determinant of cohomology.
\end{proof}

The packet
\(\texttt{certificates/hybrid/determinant\_anchor\_chi\_zero\_degeneracy}\)
verifies
\(\texttt{DETERMINANT\_ANCHOR\_CHI\_ZERO\_DEGENERACY\_VERIFIED}\):
on the branch \(\chi_\eta=0\), the determinant anchor is constant on
the full \(E\)-translation orbit, and freeness of the retained
translation action gives a rank-one separation defect for
determinant-anchor-only data.  It records that extra anchor data is
mandatory.  It does not itself construct that data, define
\(o^\lambda_R\), prove losslessness, prove quotient descent, choose
stabilizer linearizations, or prove transition compatibility.

\begin{proposition}[Extra anchor data from the degree-zero Jordan--Hoelder divisor]
\label{prop:determinant-anchor-extra-data-separates-rank-two}
\textnormal{[proved]}
Let \(V\) be a semistable vector bundle of rank \(r\) and degree \(0\)
on \(E\).  Write the associated graded object of a Jordan--Hoelder
filtration as
\[
\operatorname{gr}^{\mathrm{JH}}V=L_1\oplus\cdots\oplus L_r,
\qquad
L_i\in\operatorname{Pic}^0(E),
\]
and set
\[
\operatorname{sp}^{\mathrm{JH}}(V)
=[L_1]+\cdots+[L_r]\in
\operatorname{Sym}^r\operatorname{Pic}^0(E).
\]
This divisor is independent of the Jordan--Hoelder filtration.  Its
Abel sum is \(\det V\).  For
\[
V_0=\mathcal O_E^{\oplus2},\qquad
V_L=L\oplus L^{-1},\qquad L\ne\mathcal O_E,
\]
one has
\[
\det V_0\simeq\det V_L\simeq\mathcal O_E,\qquad
\operatorname{sp}^{\mathrm{JH}}(V_0)=2[\mathcal O_E],\qquad
\operatorname{sp}^{\mathrm{JH}}(V_L)=[L]+[L^{-1}],
\]
and the two degree-two divisors are unequal in
\(\operatorname{Sym}^2\operatorname{Pic}^0(E)\).
\end{proposition}

\begin{proof}
By Atiyah's classification and Tu's S-equivalence description for
semistable bundles on an elliptic curve
\cite{AtiyahVectorBundlesElliptic,TuSemistableElliptic}, every stable
Jordan--Hoelder factor of a degree-zero semistable vector bundle on
\(E\) is a degree-zero line bundle, and the unordered collection of
these factors is the S-equivalence class of \(V\).  Hence
\(\operatorname{sp}^{\mathrm{JH}}(V)\) is independent of the chosen
Jordan--Hoelder filtration.  The determinant of the associated graded
object is
\[
\det\operatorname{gr}^{\mathrm{JH}}(V)=L_1\otimes\cdots\otimes L_r,
\]
which equals \(\det V\); under the group law on
\(\operatorname{Pic}^0(E)\), this is precisely the Abel sum of the
divisor \(\operatorname{sp}^{\mathrm{JH}}(V)\).  The displayed
rank-two pair has determinant \(\mathcal O_E\) on both sides.  If
\(2[\mathcal O_E]=[L]+[L^{-1}]\) in
\(\operatorname{Sym}^2\operatorname{Pic}^0(E)\), then the unordered
multisets \(\{\mathcal O_E,\mathcal O_E\}\) and
\(\{L,L^{-1}\}\) are equal, so \(L\simeq\mathcal O_E\), contrary to the
hypothesis.  Thus the determinant anchor agrees while the
Jordan--Hoelder divisor separates the two objects.
\end{proof}

The packet
\(\texttt{certificates/hybrid/determinant\_anchor\_extra\_anchor\_data}\)
verifies
\(\texttt{DETERMINANT\_ANCHOR\_EXTRA\_ANCHOR\_DATA\_VERIFIED}\):
on a row where the elliptic pushforward is a degree-zero semistable
vector bundle, the additional anchor datum is the
Jordan--Hoelder S-equivalence divisor in
\(\operatorname{Sym}^r\operatorname{Pic}^0(E)\).  Its Abel sum is the
determinant anchor, and it separates the determinant-zero rank-two test
pair
\(\mathcal O_E^{\oplus2}\) and \(L\oplus L^{-1}\) with
\(L\ne\mathcal O_E\).  It does not prove global anchor losslessness,
prove \(o^\lambda_R=0\), prove quotient descent, choose stabilizer
linearizations, populate every finite wrapped prequotient, or prove
transition compatibility.

\begin{lemma}[Connected $E$-descent: separated cohomological classes]
\label{lem:connected-e-descent}
For the connected component $E^{\circ}$ of the identity in the
$E$-translation action,
\[
H^*(BE; \mathbb F_2) = \mathbb F_2[u_1, u_2],\qquad
|u_i| = 2.
\]
Hence
\[
H^1(BE; \mathbb F_2) = 0,\qquad
H^2(BE; \mathbb F_2) = \mathbb F_2 u_1 \oplus \mathbb F_2 u_2.
\]
There is no degree-one connected character; there is a degree-two
connected gerbe class
\[
\alpha^{E,\mathrm{free}} = a_1 u_1 + a_2 u_2 \in H^2(BE; \mathbb F_2).
\]
(proved, \cite{NIKULIN}.)
\end{lemma}

\begin{proof}
$E^{\mathrm{top}}\simeq T^2$ has classifying space $BE\simeq BT^2$ and
$H^*(BT^2; \mathbb F_2) = \mathbb F_2[u_1, u_2]$ with both $u_i$ of
degree $2$ (mod-$2$ Chern classes of the universal line bundles).
$H^1$ vanishes; $H^2$ is the displayed two-dimensional $\mathbb F_2$
vector space.  Vanishing of the connected gerbe class is the equation
$a_1 = a_2 = 0$ on the actual Borel/edge class of the equivariant
reduced determinant complex, not a consequence of $H^1(BE) = 0$ or
ordinary translation invariance.
\end{proof}

\begin{lemma}[{Finite stabilizer $E[N]$-descent: gerbe and linearization classes}]
\label{lem:finite-stabilizer-e-descent}
For the order-two finite subgroup $E[2]\cong(\Z/2)^2$,
\[
H^*(BE[2]; \mathbb F_2) = \mathbb F_2[x_1, x_2],\qquad
|x_i| = 1.
\]
The degree-two gerbe class is
\[
\beta^{E, E[2]}
= b_{20} x_1^2 + b_{11} x_1 x_2 + b_{02} x_2^2,\qquad
(b_{20}, b_{11}, b_{02}) = (r_1,\,r_1 + r_2 + r_3,\,r_2),
\]
detected on cyclic subgroup restrictions by joint reduction; the $r_i
\in\mathbb F_2$ are the three nonzero point-stabilizer parities.
After $\beta^{E, E[2]} = 0$ (i.e.\ $r_1 = r_2 = r_3 = 0$), a separate
degree-one linearization class remains:
\[
\lambda^{E, E[2]} = \lambda_1 x_1 + \lambda_2 x_2 \in H^1(BE[2]; \mathbb F_2),
\qquad \rho_1 = \lambda_1,\ \rho_2 = \lambda_2,\ \rho_3 = \lambda_1 + \lambda_2.
\]
The $r_i$ are gerbe bits in $H^2$; the $\rho_i$ are linearization bits
in $H^1$; the two live in distinct cohomological degrees and remain
independent.
For even $N\ge 4$ with $2^a\parallel N,\ a\ge 2$,
\[
H^*\!\bigl(B(\Z/2^a)^2; \mathbb F_2\bigr)
= \Lambda(x_1, x_2)\otimes\mathbb F_2[y_1, y_2],
\]
\[
\beta^{E, N} = A_1^{(N)} y_1 + A_{12}^{(N)} x_1 x_2 + A_2^{(N)} y_2,
\qquad
\lambda^{E, N} = \lambda_1^{(N)} x_1 + \lambda_2^{(N)} x_2.
\]
The mixed term $A_{12}^{(N)}$ is invisible to cyclic order-two restrictions;
even $N\ge 4$ tests must be checked on the full two-primary generators.
For odd $N$, the mod-2 orientation obstruction vanishes by transfer.
(proved.)
\end{lemma}

\begin{proof}
$E[2]\cong(\Z/2)^2$ has $H^*(B(\Z/2)^2; \mathbb F_2) = \mathbb F_2[x_1,
x_2]$ with $|x_i| = 1$.  The class
\(\beta^{E,E[2]}\) is the finite-stabilizer Cech--Borel edge class,
not the restriction of the connected \(BE\)-class
\(\alpha^{E,\mathrm{free}}\).  Its restriction to the three non-zero
cyclic subgroups gives the three parities \(r_i\); inversion of this
three-by-three system gives
\((b_{20},b_{11},b_{02})=(r_1,r_1+r_2+r_3,r_2)\).
The degree-one linearization
class $\lambda^{E, E[2]}$ is the residual character of the underlying
line bundle after $\beta = 0$ is imposed; it lives in $H^1$ and is
distinct from the gerbe $\beta\in H^2$ in degree.  Even $N\ge 4$
introduces the polynomial generators $y_i \in H^2(B(\Z/2^a))$ and the
mixed term $x_1 x_2$, which on cyclic restrictions vanishes; full
two-primary detection is needed.  Odd $N$ contributes only $\mathbb
F_p$-cohomology with $p$ odd, on which the mod-$2$ orientation has zero
restriction by the transfer.
\end{proof}

\begin{remark}[Translation linearization criterion, not vanishing]
\label{rem:translation-linearization-criterion}
For a stabilizer $H \subset E[N]$ and a determinant anchor line bundle
$\lambda(\mathcal F)$, ordinary translation invariance $g^*\lambda
\simeq \lambda$ in $\mathrm{Pic}^0(E)$ does not choose coherent
isomorphisms $\phi_g : \lambda \to g^*\lambda$.  After the square-root
gerbe is killed, the choices form a character torsor in $H^1(BH;
\mathbb F_2)$, and quotient orientation requires the linearization
class
\[
\lambda^H \in H^1(BH; \mathbb F_2)
\]
to vanish.  This is a \emph{criterion} for descent on the wrapped
stratum; it is not implied by $g^*[\lambda] = [\lambda]$ in
$\mathrm{Pic}(E)$.  Proposition~\ref{prop:determinant-anchor-translation-weight}
computes the translation law for \(\lambda^{\det}\) and its weight
\(\chi(\mathcal F)\).  The finite-cover normalisation to unit weight
when \(\chi(\mathcal F)\ne0\) and the coherent stabilizer
linearization are later anchor rows.  The \(\chi(\mathcal F)=0\)
determinant-anchor degeneracy is
Proposition~\ref{prop:determinant-anchor-chi-zero-degeneracy}; the
degree-zero Jordan--Hoelder divisor of
Proposition~\ref{prop:determinant-anchor-extra-data-separates-rank-two}
separates the determinant-zero rank-two test pair.  Full anchor
losslessness, finite-stage population, quotient descent, transition
compatibility, and the coherent stabilizer linearization remain later
rows.  Until those rows are supplied, this paragraph is only the
descent criterion for a chosen linearized anchor.
(proved as a criterion; determinant-anchor weight proved in
Proposition~\ref{prop:determinant-anchor-translation-weight}.)
\end{remark}

\begin{remark}[Mod-2 character does not control odd-order anchor
characters]
\label{rem:mod2-not-odd-order}
The mod-$2$ orientation character does not by itself control odd-order
anchor characters.  For anchor descent on the wrapped stratum, also
compute the actual stabilizer character on the chosen anchor
trivialization, valued in $\widehat H = \mathrm{Hom}(H, \C^\times)$, or
define the retained stabilizer as preserving that trivialization.
(unverified; supplied as part of the finite Hall datum.)
\end{remark}

\subsection{Cusp polarization and chamber transport}
\label{subsec:cusp-polarization-chamber}

\begin{definition}[Cusp polarization on $\Gamma_{\mathrm{eff}}$]
\label{def:cusp-polarization}
The chamber convention $0 < |s|\ll|q|\ll|r|^{-1}\ll 1$ at the standard
cusp $\mathfrak c_\infty$ determines the cusp-positive semigroup
$\Gamma_{\mathrm{eff}}$ of Definition~\ref{def:gram-index-lattice}.  The
order on triples is lexicographic: $m > 0$ first because $s$ is
smallest; if $m = 0$ then $n > 0$ because $q$ is next; if $m = n = 0$
then $l < 0$ because $|r|^{-1}\ll 1$.  This is a chamber convention for
the completed product, not a microscopic charge positivity statement.
\end{definition}

\begin{remark}[Chamber transport]
\label{rem:chamber-transport}
The pairing $\langle\gamma, \gamma'\rangle_{\mathrm{ind}} = 2(n m' + n'
m) - l l'$ is the Gram pairing on triples; the order on
$\Gamma_{\mathrm{eff}}$ is a Fock polarization at one cusp.  An element
of the duality group $\Sp_4(\Z)$ transports both the chamber and the
expansion to the corresponding product chamber at the transported cusp.
The arithmetic group $\Sp_4(\Z)$ is the symmetry of the genus-two
chemical potential $Z$ and of the invariant Gram plane on
$\Lambda^{3, 2}$; it is not asserted to be the full microscopic
T-/U-duality group of the K3$\times T^2$ compactification, which is
larger (folklore, \cite{HM,DVV}).
\end{remark}

\begin{remark}[Chamber dependence of the wrapped stratum]
\label{rem:chamber-dependence-wrapped}
The cusp chamber chooses a positive direction on
$\Gamma_{\mathrm{eff}}$.  At $\mathfrak c_\infty$ the wrapped stratum
carries $\deg_E > 0$ branes; at a transported cusp the wrapped
direction is the image of $\deg_E$ under the corresponding chamber
transport.  Mixed Hall correspondences are chamber-equivariant when the
finite Hall stage $R$ is compatible with the duality element; otherwise
the correspondence stack is the original closed flag-Quot stack
restricted to the new chamber.
\end{remark}

\begin{definition}[Chamber-compatible Hall descent datum]
\label{def:chamber-compatible-hall-descent}
Let \(\sigma\) and \(\sigma'\) be two chosen stability data, and let
\(g\in\Sp_4(\Z)\) transport the standard cusp chamber to another
product chamber.  A chamber-compatible Hall descent datum is a
Kontsevich--Soibelman wall-crossing system
\[
\mathsf{WC}_{g,R}:\mathcal C_{\sigma,S,\le R}\longrightarrow
\mathcal C_{\sigma',S,\le R'}
\]
on the retained finite Hall stages, with transition maps in \(R\),
which preserves the normal-ordered Gram map \(\overline\Pi_X\), transports
the local/wrapped split, the quotient-after-correspondence datum, the
orientation lines, and the protected integration maps.  The identity
\(\overline\Pi_X(\widehat c\star\widehat c')
=\overline\Pi_X(\widehat c)+\overline\Pi_X(\widehat c')\) on
\(\widehat\Gamma_X\) is a chamber-independent lattice fact; the compact
Hall category and its primitive bracket are not chamber-independent without
\(\mathsf{WC}_{g,R}\) (folklore/\(\beta\), \cite{KSCoHA}).
\end{definition}

\section{Local protected observable algebra (sectorial source datum)}
\label{sec:local-protected-observable-algebra}

The compact realization route of Chapter~\ref{chap:dirac-igusa-pro-object}
takes as supplied input a chain-level local observable system on
opens $U\subset X = S\times E$.  The system is graded by the
algebraic/effective Hall support, factorizes on disjoint opens, and
restricts on $E$-projection-finite supports to a chiral object on
$\operatorname{Ran}(E)_{b = 0}$.  The data are external assumptions
named in the literature and form the source datum for the compact
realization.

\begin{definition}[Local protected observable algebra]
\label{def:local-protected-observable-algebra}
Let $U\subset X = S\times E$ be open.  Let
$\Gamma_{X, \sigma}^{\mathrm{eff}}(U)\subset\Gamma_X^{\mathrm{alg}}
\subset\Gamma_X^{\mathrm{form}}$ be the supplied algebraic/effective
charge support.  For $c\in\Gamma_{X, \sigma}^{\mathrm{eff}}(U)$, let
$\mathfrak M_c(U)$ be the supplied derived moduli stack of compactly
supported B-branes on $U$ of Mukai charge $c$.  The PTVV theorem supplies
the ambient $(-1)$-shifted symplectic structure on the $\mathrm{CY}_3$
perfect-complex moduli problem (proved, \cite{PTVV2013}).  Let
$\Phi_c$ denote the BBDJS perverse sheaf $P^\bullet_{\mathfrak M_c, s_c}$
on an oriented d-critical truncation, when such a d-critical structure
and square-root orientation are supplied (proved,
\cite[Theorem~6.9]{BBDJS2015}).  Let $o_c$ be the orientation line, a
square root of the determinant line of the cotangent complex of
$\mathfrak M_c(U)$, supplied as part of the Hall datum.

The \emph{local protected observable algebra of $X$ on $U$} is the
source object
\[
\mathcal A_X(U)
:= \bigoplus_{c\in\Gamma_{X, \sigma}^{\mathrm{eff}}(U)}
H_*^{\mathrm{BM}}\!\bigl(\mathfrak M_c(U),\,\Phi_c\otimes o_c\bigr),
\]
graded by the supplied algebraic/effective support, not by the full
formal lattice $\Gamma_X^{\mathrm{form}}$.  For pairwise disjoint
$U_1, \ldots, U_k\subset X$, the proposed factorization map is
\[
\mathcal A_X(U_1)\otimes\cdots\otimes\mathcal A_X(U_k)
\longrightarrow\mathcal A_X(U_1\sqcup\cdots\sqcup U_k),
\]
defined chain-level by Massey--Saito Thom--Sebastiani for vanishing
cycles
\[
\Phi_{F_1\oplus\cdots\oplus F_k}
\simeq \Phi_{F_1}\boxtimes\cdots\boxtimes\Phi_{F_k}
\]
on disjoint d-critical charts (proved,
\cite{MasseyST,SaitoMHM}), together with the Ext-vanishing
$\mathrm{Ext}^*(\mathcal F, \mathcal G) = 0$ for branes
$\mathcal F, \mathcal G$ with disjoint supports
(folklore).
\end{definition}

\begin{remark}[External source datum]
\label{rem:source-interface-external-assumptions}
The components are external assumptions: PTVV $(-1)$-shifted symplectic
structure on derived moduli of B-branes (proved,
\cite{PTVV2013}); BBDJS vanishing-cycle perverse sheaf and its
Thom--Sebastiani isomorphism (proved,
\cite{MasseyST,SaitoMHM,BBDJS2015}); orientation lines and
Joyce--Upmeier-compatible multiplicative orientation data (ambient
technology proved in \cite[Theorem~4.1]{JoyceUpmeier}; the
\(K3\times E\) quotient orientation is supplied separately).  The factorization map relies on disjoint-supports
Ext-vanishing for compactly supported branes (folklore;
absent for non-compact supports, where one recovers chiral algebra
structure, not commutative factorization).  Compactness on $X = S\times
E$ is open because the relevant moduli of branes are typically
non-quasicompact in $\Lambda_S$-charge; the sectorial truncation of
\S\ref{sec:k3xe-hybrid-ran-carrier} is the finite-stage carrier, and
the compact-source recognition datum \((S1)\)--\((S10)\) is the
remaining finite recognition problem.
At finite stage the retained-moduli assertion is represented by the
tables \(\mathrm{hn\_type\_bounds}\),
\(\mathrm{semistable\_substacks}\), \(\mathrm{derived\_enhancements}\),
\(\mathrm{universal\_complexes}\), \(\mathrm{scalar\_rigidifications}\),
\(\mathrm{stratifications}\), \(\mathrm{extension\_flag\_stacks}\),
\(\mathrm{cosection\_atlas}\), \(\mathrm{transitions}\), and
\(\mathrm{scalar\_firewall}\), not by the scalar Hilbert-scheme test.
At present the packet
\(\texttt{certificates/moduli/k3e\_finite\_moduli}\) supplies only the
obstruction ledger
\(\texttt{MODULI\_OBSTRUCTION\_LEDGER\_VERIFIED}\), with
\(\texttt{moduli\_certification=false}\).
\end{remark}

\section{The four obstructions and the comparison target}
\label{sec:k3xe-four-obstructions}

The hybrid carrier and the source datum are supplied as a finite
holomorphic $E_3$-prefactorization datum and a \(K3\)-to-\(E\)
specialization datum on $\operatorname{Ran}^{\mathrm{hyb}}(E)$.  The
cross-level comparison
$\operatorname{Pf}_{\mathrm{prot}}(\mathfrak D_X^{\mathrm{DI}}) =
\Delta_5$ holds on \(K3\times E\) relative to four retained data,
with \((O2)\) represented by the retained wall-atlas datum of
Appendix~\ref{app:obstruction-discharges}
(Theorems~\ref{thm:G-D0}, \ref{thm:G-O1}, \ref{thm:G-O1-plus},
\ref{thm:G-O2-orbit-transport}).  They are the orientation and
Pfaffian data feeding
Chapter~\ref{chap:dirac-igusa-pro-object}.

\begin{description}
\item[$(\mathrm{D0})$ {\rm D0-degeneration limit.}]
The retained D0-HN datum of Definition~\ref{def:G-D0-HN-system}
supplies the flat D0-degeneration family, additive K3 semiregularity
cosections, reduced vanishing-cycle and orientation specializations,
Pfaffian-line extension, and strict HN transitions
(criterion Theorem~\ref{thm:G-D0}).
\item[$(\mathrm{O1})$ {\rm Strong reduced orientation.}]
The K3$\times$E-quotient self-Ext frame bundle carries a strong reduced
orientation on the wrapped stratum, descending the BBDJS orientation
line $o_c$ across the $E$-quotient relative to
\((O1_{\mathrm{quot}})\) (Theorem~\ref{thm:G-O1}).
\item[$(\mathrm{O1}^+)$ {\rm Weyl-equivariant transport.}]
The orientation transports $W^{(2)}(\Lambda_{II}^{2, 1})$-equivariantly
along reflections relative to the chosen Weyl lifts
(Theorem~\ref{thm:G-O1-plus}); after the type-II wall atlas supplies
the rank-one local signs, the transported character agrees with
$\nu_{\Delta_5}$ on the BKM Weyl group.
\item[$(\mathrm{O2})$ {\rm Local Pfaffian wall normal form.}]
On type-II walls of $\Lambda_{II}^{2, 1}$, the local Pfaffian wall
normal form is the rank-one crossing
$[\mathbb R\xrightarrow{u}\mathbb R]$; the retained
\((O2_{\mathrm{atlas}})\) datum transports these local signs over the
half-Hilbert orbit, whose sign support has cardinality
$4096 = 2^{12}$ (Theorem~\ref{thm:G-O2-orbit-transport}).
\end{description}

The protected scalar shadow on the closed K3$\times$E$\to$point
cobordism is the OP/DT identity
\[
Z^X_{\mathrm{DT}}
= Z^X_{\mathrm{OP}}(P_{\mathrm{OP}}, Q_{\mathrm{OP}}, T_{\mathrm{OP}})
= -4096\,\Delta_5(Z)^{-2}
\qquad\text{(proved, \cite{OB,OPand,DT}).}
\]
Equivalently $\chi_{10}^{\mathrm{OP}} = (64^{-1}\Delta_5)^2 = 4096^{-1}\Delta_5^2$
and $Z^X_{\mathrm{OP}} = -(\chi_{10}^{\mathrm{OP}})^{-1}$; the factor
$4096 = 64^2$ is a theta-leading normalization conversion and the minus
sign is the OP scalar convention.  The orientation-forgetful shadow
\[
Z_{\mathrm{BPS}}^{K3\times E}
= (\Phi_{10}^{\mathrm{un}})^{-1}
= \Delta_5^{-2}
\]
is the level-$n$ protected trace of the retained \(K3\times E\) data on the
closed cobordism in the unnormalized convention $\Phi_{10}^{\mathrm{un}}
= \Delta_5^2$; the inverse-square scalar cannot recover the orientation
character that distinguishes $\Delta_5$ from $-\Delta_5$.  A
\(3d\) gravitational path-integral realization requires the Hall--Borcherds residual
of \cite{ChiralBarCobarVolII} and is not claimed.  The conditional
cross-level identity
$\operatorname{Pf}_{\mathrm{prot}}(\mathfrak D_X^{\mathrm{DI}}) =
\Delta_5$ supplies the orientation/Pfaffian datum under
$(\mathrm{D0}), (\mathrm{O1}), (\mathrm{O1}^+), (\mathrm{O2})$.

\begin{remark}[Compatibility of the \(\kappa\)-tuple]
\label{rem:kappa-tuple-compatibility}
The $\kappa$-tuple $(0, 0, 3, 5, 24)$ is consistent across
\cite{CYQGVolIII} (Vol III), the Hall--Borcherds
residual of \cite{ChiralBarCobarVolII} (Vol II), and the
Beilinson--Drinfeld factorization framework of \cite{BD} via Vol I
\cite{ChiralBarCobarVolI}.  The local model
$\R^2_{\mathrm{top}}\times\C^2_{\mathrm{hol}}$ + brane completion of
\cite{MixedHTSCompanion} contributes the
holomorphic de Rham obstruction class on the global target; the
identity $Z_{\mathrm{BPS}}^{K3\times E} = \Delta_5^{-2}$ does not
promote to the global target via formal Hamiltonian BF, and the
Hall--Borcherds residual is the remaining chain-level bridge.
\end{remark}

\section{Notation and fixed data}
\label{sec:k3xe-notation-summary}

\begin{description}
\item[$X = S\times E$.]
Smooth complex projective trivial-canonical threefold; $S$ is a K3
surface, $E$ a smooth elliptic curve.  $\dim_\C X = 3$ and
\(h^{1,0}(X)=1\).
\item[$\Lambda_S$.]
Mukai lattice; even unimodular signature $(4, 20)$.
$\Lambda_S\cong U^{\oplus 4}\oplus E_8(-1)^{\oplus 2}\cong II_{4, 20}$.
\item[$\Gamma_X^{\mathrm{form}}, \Gamma_X^{\mathrm{phys}}$.]
Formal full-Mukai dyonic lattice $\Lambda_S\oplus\Lambda_S$.  The
superscript ``$\mathrm{phys}$'' denotes this formal even branch, not the
$\mathcal N = 4$ charge lattice.
\item[$\Gamma_X^{\mathrm{alg}},\Gamma_{X, \sigma}^{\mathrm{eff}}$.]
Algebraic K-class lattice $\subset\Gamma_X^{\mathrm{form}}$; chosen
$\sigma$-effective semigroup $\subset\Gamma_X^{\mathrm{alg}}$.
\item[$\Gamma_{\mathrm{gram}} = \Z^3$.]
Gram-index lattice, $(n, l, m)$.
\item[$\Pi_X$.]
Quadratic Gram map, $\Pi_X(Q, P) = ((Q, Q)/2,\,(Q, P),\,(P, P)/2)$.
\item[$B$.]
Symmetric bilinear cocycle, $B(c, c') = (Q\cdot Q',\,Q\cdot P' + Q'\cdot
P,\,P\cdot P')$; satisfies $B = -\delta\Pi_X$ and $B(c, c) = 2\Pi_X(c)$.
\item[$\widehat\Gamma_X$.]
Normal-ordered extension $\Gamma_X^{\mathrm{form}}\oplus_B\Gamma_{\mathrm
{gram}}$; abelian under $\star$.
\item[$\overline\Pi_X$.]
Additive normal-ordered Gram homomorphism $\overline\Pi_X(c, T) =
\Pi_X(c) - T$.
\item[$\operatorname{Ran}^{\mathrm{hyb}}(E)$.]
Finite-stage hybrid prestack on the elliptic factor, with local
\(b=0\) subprestack
\(\operatorname{Ran}^{\mathrm{loc}}_{R,\mathrm{pre}}(E)\) and wrapped
\(b>0\) subprestack
\(\operatorname{Ran}^{\mathrm{wr}}_{R,\mathrm{pre}}(E)\).
\item[$\mathcal D_X = \Delta_5$, $\mathrm{den}(\mathfrak g_{\Delta_5})$,
$\mathrm{Pf}_{\mathrm{prot}}(\mathfrak D_X^{\mathrm{DI}})$.]
Three named weight-five objects: automorphic section, BKM denominator,
compact Pfaffian (View \textcircled{1}, \textcircled{2}, \textcircled{3}).
\end{description}

The fixed data are: (a) the formal Mukai sector
$\Gamma_X^{\mathrm{form}} = \Lambda_S\oplus\Lambda_S$ and its three-tier
support hierarchy; (b) the bilinear polarization \(B = -\delta\Pi_X\) as the
full structural content of the Gram defect; (c) the normal-ordered
extension $\widehat\Gamma_X$ with additive homomorphism
$\overline\Pi_X$, fixed by the fixed-lift raw-pushforward obstruction
theorem; (d) the hybrid Ran carrier
$\operatorname{Ran}^{\mathrm{hyb}}(E)$, fixed by the projection-locality
obstruction; (e) the connected/finite split of the $E$-descent classes,
with the linearization criterion separated from the gerbe class; (f)
the local protected observable algebra source datum, taken as an
external assumption; (g) the four obstructions
$(\mathrm{D0}), (\mathrm{O1}), (\mathrm{O1}^+), (\mathrm{O2})$
conditioning the cross-level identity ②$\leftrightarrow$③.

Chapter~\ref{chap:dirac-igusa-pro-object} constructs from these fixed
data the Dirac--Igusa pro-object $\mathfrak D_X^{
\mathrm{DI}} = \lim_R(\mathcal A_{X, R}^{E_3}, F_{X, R}^{\mathrm{hyb}},
\Gamma_{X, R}, \Pi_X, \Phi_R, o_R, H_R, P_R^\Pi, C_{X, R},
\Theta_{\mathrm{Kos}, R}, \mathcal L_{\mathrm{Pf}, R}, \mathrm{pf}_{X, R},
\epsilon_{o, R}, \mathrm{Rec}_R)$.
